% SPDX-FileCopyrightText: 2026 Will Cook
% SPDX-License-Identifier: CC-BY-4.0
%
% Long-form reasoning record seeded from the complete problem note.
% The short note and this record are now independent publication units.
% This flat manuscript is generated from reviewable parts below.
\documentclass[11pt]{article}
\input{problem-note-preamble}
\usepackage{longtable}
\usepackage{colortbl}
\usepackage[htt]{hyphenat}
\usepackage{newunicodechar}
\newunicodechar{ℕ}{\ensuremath{\N}}
\newunicodechar{ℚ}{\ensuremath{\Q}}
\newunicodechar{ℤ}{\ensuremath{\Z}}
\newunicodechar{₀}{\ensuremath{{}_0}}
% Problem-specific source pin: #269's local-window consumer advances
% independently of the corpus default.
\renewcommand{\commit}{ee650b32b8b2cb98b94e5500df5370d85f7403b8}
\renewcommand{\commitshort}{ee650b32b8b2}
\hypersetup{
  pdftitle={The three-prime running least common multiple: complete reasoning record for Erdős Problem 269},
  pdfsubject={Erdős Problem Notes: Problem 269},
  pdfkeywords={transcendence, Hecke-Mahler values, least common multiple, smooth numbers, lattice sums, Lean 4},
}

% The three pure-power coordinates and the running height.
\newcommand{\hgt}{\operatorname{H}}
\newcommand{\lo}{\lambda}
\newcommand{\hi}{\eta}
\newcommand{\Lprefix}{\operatorname{L}}
\newcommand{\Kernel}{\operatorname{K}}
\newcommand{\lpr}{\operatorname{lpr}}
\newcommand{\Esc}{\mathsf E}
\newcommand{\authnone}{\textcolor{BrickRed!75!black}{\sffamily\bfseries No result here}}
\newcommand{\authliterature}{\textcolor{BurntOrange!85!black}{\sffamily\bfseries Literature}}
\newcommand{\authlean}{\textcolor{MidnightBlue!82!black}{\sffamily\bfseries Paper + Lean}}
\newcommand{\authexternal}{\textcolor{Violet!80!black}{\sffamily\bfseries Paper + external theorem}}
\newcommand{\authcompute}{\textcolor{Violet!80!black}{\sffamily\bfseries Direct computation}}

\runninghead{Erd\H{o}s \#269: complete reasoning record}
\title{The Three-Prime Running LCM: Complete Reasoning Record}
\subtitle{Erd\H{o}s Problem~\#269}
\author{Will Cook}
\date{July 2026}

\setlength{\emergencystretch}{2em}
\def\UrlBreaks{\do\.\do\_\do\/}
\begin{document}
% ---- part core ----
\maketitle
\noteprovenance

\begin{abstract}
\noindent
For every pair of distinct primes \(p,q\) both series of Erd\H{o}s
Problem~\#269 are transcendental: the original series \(\mathcal R_{p,q}\), in
which a reciprocal of the running least common multiple is taken at every
\(\{p,q\}\)-smooth integer, and the de-duplicated series \(\mathcal D_{p,q}\),
in which each distinct running least common multiple contributes once.  The
hypotheses are that \(p\) and \(q\) are prime and distinct, and nothing else.
Every instance of the problem with \(|P|=2\) is therefore settled, at the
stronger level of transcendence.  Ordering the two pure-power channels gives a
Beatty word, both values become explicit polynomials in one Hecke--Mahler
boundary sum, and transcendence of that sum is the theorem of Loxton and van
der Poorten in the modern form of Bugeaud and Laurent.  This is an argument in
this paper over a cited external theorem, and no Lean declaration formalises
it.  Steve Fan posted the same factorisation, reduction and conclusion first,
on 26 June 2026, so the priority for the two-prime case is his.

The third prime changes the linear algebra of the problem, and the change is
exact.  Write \(\Kernel(i,j,k)=1/\Lprefix(p^{i}q^{j}r^{k})\) for the reciprocal
running-LCM kernel.  At two generators the kernel is an outer product and every
two-by-two minor vanishes.  At three pairwise distinct generators, for every
order \(n\) there are injective index families \(I,J\) whose minor
\(\det(\Kernel(I(a),J(b),k))_{0\le a,b<n}\) is nonzero simultaneously in every
layer \(k\), so for no finite \(d\) do there exist rational-valued
\(f_\ell(i)\) and \(G_\ell(j,k)\) with
\(\Kernel(i,j,k)=\sum_{\ell<d}f_\ell(i)G_\ell(j,k)\).  Both halves of that
transition are checked by the Lean kernel at the pinned source
revision~\commitshort.  A single binary floor carry is the whole difference,
and after row and column normalisation that carry matrix admits no better
uniform approximation by a finite separated sum than \((r-1)/(2r)\), which is
half its own jump and is attained by a constant.  These statements exclude
classes of representations and settle no instance of the problem.

For the repeated \(\{2,3,5\}\) series the note constructs the literal infinite
object.  The dyadic shell masses are summable, their normalised tails \(X_a\)
satisfy the integer-coefficient recurrence \(X_{a+1}=b_aX_a-m_a\) with
\(b_a\in\{2,6,10,30\}\), and \(0<X_a\le(n_a^{2}+8n_a+18)/9\) where \(n_a\) is
the total height rank at \(2^{a}\).  A rational value with denominator
\(2^{u}3^{v}5^{w}B\), \(\gcd(B,30)=1\), forces \(BX_a\) to be a positive
integer from the explicit onset \(a_D=u+1+2v+3w\), bounded by
\(K(B,a)=\lfloor B(j_a^{2}+10j_a+27)/9\rfloor\) at the endpoint jump index
\(j_a=n_a-1\).  Let \(\Esc(G)\) be cofinal local-window residue escape against
a bound \(G\).  Then \(\Esc(G)\) is equivalent to irrationality of the series
for every \(G\) with \(K\le G\) and \(G(B,a)=o(2^{a})\), and the lower half of
that band cannot be removed, since \(\Esc(0)\) is vacuous.  Nonintegrality of
every reduced tail is equivalent to irrationality as well.  Both remaining
questions are therefore restatements of the problem, and neither of them is a
weaker sufficient condition.

The finite evidence is exact and bounded.  The pinned integer checker finds an
escaping window of length at most \(18\) for each of \(3{,}869{,}934\)
denominator/start pairs, covering every \(B\le5000\) coprime to \(30\) and
every start \(100\le\ell\le3000\); a window-growth law
\(8^{h}/15<W_{\ell,h}<15\cdot8^{h}\) shows that this behaviour cannot persist
at bounded length for all starts, so the scan measures a bounded region and
leaves the cofinal quantifier untouched.  A window-\(128\) lattice certificate at launch
\(a_1=10005\) excludes an arithmetically described family of rational values,
and a continued-fraction certificate excludes every reduced denominator below
\(2^{22482}\).  Erd\H{o}s~\#269 remains open for every \(|P|\ge3\).
\end{abstract}

\section{The problem, and what is settled}\label{long269:sec:problem}

Let $P$ be a finite set of primes with $|P|\ge2$, and let
$a_1<a_2<\cdots$ enumerate the positive integers all of whose prime factors lie
in $P$.  Erd\H{o}s Problem~\#269 asks whether
\[
 \sum_{n\ge1}\frac{1}{[a_1,\ldots,a_n]}
\]
is irrational, where $[a_1,\ldots,a_n]$ is the least common
multiple~\cite[p.~65]{erdosgraham1980}\cite[p.~106]{erdos1988}.  Bloom's
catalogue records the problem as open~\cite{erdosproblems}.  This note settles
every instance with $|P|=2$, at the level of transcendence, and the unresolved
finite cases begin at $|P|=3$.

The restriction $|P|\ge2$ is necessary.  For $P=\{p\}$ the enumeration is
$a_n=p^{\,n-1}$, so $[a_1,\ldots,a_n]=p^{\,n-1}$ and the sum is $p/(p-1)$.  For
infinite $P$ the sum is always irrational, which Erd\H{o}s calls a simple
exercise~\cite[p.~106]{erdos1988}.

Write $\mathcal R_P$ for the sum above and $\mathcal D_P$ for the
de-duplicated sum, in which each distinct value of the running least common
multiple contributes its reciprocal once.  The two differ because the running
value is constant along stretches of the enumeration.  In a letter written on
1~January 1973 Erd\H{o}s recorded that he could prove irrationality of the
de-duplicated sum~\cite[p.~335]{erdos1974letter}.  He states that assertion for
a general finite list of primes and supplies no proof, so the historical record
already covers $\mathcal D_P$ for every finite $P$; the open question addressed
by the catalogue is $\mathcal R_P$.

\begin{theorem}[two-prime transcendence]\label{long269:res:lead-two-prime}
Let $p$ and $q$ be distinct primes.  Then $\mathcal R_{\{p,q\}}$ and
$\mathcal D_{\{p,q\}}$ are transcendental.
\end{theorem}

Section~\ref{long269:sec:two-prime} proves this.  Both values are explicit polynomials
in one Hecke--Mahler boundary sum $A$, the coefficients are nonzero rationals,
and transcendence of $A$ is the theorem of Loxton and van der Poorten
\cite[p.~40, Theorem~8]{loxtonvdp1977} in the modern form of Bugeaud and Laurent's
Theorem~1.1 \cite[p.~61, Theorem~1.1]{bugeaudlaurent2023}.  The evidence class
is an ordinary proof in this paper over a cited external theorem, and no Lean
declaration formalises it.

\paragraph{Attribution.}
Steve Fan posted the running-LCM identity for every $|P|$, the two-channel
factorisation, the Hecke--Mahler reduction and the transcendence conclusion in
the discussion thread of the problem's page on 26 June
2026~\cite{fan2026comment}, with follow-up remarks extending the argument to
arbitrary coprime pairs.  The priority for the two-prime case is his.  The
proof below is independent and is given in full because it fixes the $|P|=2$
boundary of the problem and because the same boundary sum drives the
three-prime discussion.  The identity of Theorem~\ref{long269:res:lcm} is cited to him.

\paragraph{What the third prime does.}
Section~\ref{long269:sec:rank} isolates the obstruction exactly.  After separating row
and column factors, the two-prime kernel is an outer product and the
three-prime kernel retains one binary floor carry.  That carry is enough to
produce nonsingular minors of every order, uniformly in the third coordinate,
and to make the normalised carry matrix inapproximable by finite separated sums
below half its own jump.  Sections~\ref{long269:sec:blocks} to~\ref{long269:sec:escape} take a
second route on the literal $\{2,3,5\}$ series: dyadic shells give an
integer-coefficient tail recurrence, a strict endpoint clears the smooth part
of a hypothetical denominator, and the surviving carry is trapped between a
quadratic bound and a residue condition.  Section~\ref{long269:sec:evidence} reports
the finite evidence and its ceiling, and Section~\ref{long269:sec:open} states what
remains.

Throughout, $p,q,r$ are pairwise distinct primes.  Call $n$ \emph{smooth} when
$n=p^{i}q^{j}r^{k}$ for some $i,j,k\ge0$; this is the
\lword{Erdos269/ThreePrimeRunningLcm.lean}{32}{smooth3Val}{smooth lattice
value}.  For $x\ge1$ write
\[
 \Lprefix(x)=\lcm\{\,n\le x:\ n\ \text{smooth}\,\},
 \qquad
 \hgt(x)=p^{\lfloor\log_p x\rfloor}\,q^{\lfloor\log_q x\rfloor}\,
 r^{\lfloor\log_r x\rfloor},
\]
the
\lword{Erdos269/ThreePrimeRunningLcm.lean}{54}{smoothPrefixLcm}{running least
common multiple} and the
\lword{Erdos269/ThreePrimeRunningLcm.lean}{37}{threePrimeHeight}{pure-power
height}, where $\lfloor\log_b x\rfloor$ is the largest $e$ with $b^{e}\le x$.
Since $a_1,\ldots,a_n$ are exactly the smooth numbers up to $a_n$, we have
$[a_1,\ldots,a_n]=\Lprefix(a_n)$.  The reciprocal of the height at a smooth
point is the
\lword{Erdos269/ThreePrimeRunningLcm.lean}{41}{threePrimeKernelQ}{lattice
kernel}
\[
 \Kernel(i,j,k)=\frac{1}{\hgt(p^{i}q^{j}r^{k})} .
\]
Here ``smooth'' always means supported on the fixed prime set, and it is not
the varying-bound notion counted by $\Psi(x,y)$ in the Dickman--Hildebrand
theory \cite[Theorem~1]{hildebrand1986}; no smooth-number density asymptotic is
used below.  For $b\in\{p,q,r\}$ the set of positive powers of $b$ is the
\emph{$b$-channel}.

The statement of Problem~\#269 has been formalised before, as a conjecture with
an unfilled proof, in the \emph{Formal Conjectures}
collection~\cite{formalconjectures269}.  That is a formal statement of the
question up to an initial constant: its Nat-indexed series includes the
empty-prefix term, so its value is exactly \(1\) plus the conventional series.
Its rational, irrational and infinite-prime
assertions all end in \texttt{sorry}.  Kova\v{c} and Tao~\cite{kovactao2024}
treat several irrationality problems of Erd\H{o}s for series of unit fractions
by elementary means; nothing from that work is used here.

\medskip
\noindent\textbf{Keywords.} irrationality; transcendence; least common
multiple; smooth numbers; separated rank; Lean~4.
\qquad
\textbf{MSC 2020.} 11J72 (primary); 11A05, 11N25, 68V20 (secondary).

\section{The finite geometry of the running value}\label{long269:sec:lcm}

The smooth numbers up to $x$ are indexed by the exponent triples $(i,j,k)$ with
$i\le\lfloor\log_p x\rfloor$, $j\le\lfloor\log_q x\rfloor$,
$k\le\lfloor\log_r x\rfloor$ and $p^{i}q^{j}r^{k}\le x$, the
\lword{Erdos269/ThreePrimeRunningLcm.lean}{47}{smoothPrefixExponents}{smooth
prefix index set}.  Both conditions matter: the coordinate box is strictly
larger than the prefix, since a product of three large pure powers can exceed
$x$ while each factor does not.

\begin{theorem}[the running least common multiple]\label{long269:res:lcm}
Let $p,q,r$ be pairwise distinct primes and $x\ge1$.  Then
$\Lprefix(x)=\hgt(x)$.
\end{theorem}

\begin{proof}
Every smooth $n\le x$ has exponents bounded by the corresponding integer
logarithms, so $n\mid\hgt(x)$ and hence $\Lprefix(x)\mid\hgt(x)$.  Conversely
the three pure powers $p^{\lfloor\log_p x\rfloor}$,
$q^{\lfloor\log_q x\rfloor}$ and $r^{\lfloor\log_r x\rfloor}$ are themselves
smooth numbers not exceeding $x$, so each divides $\Lprefix(x)$, and distinct
primes have coprime powers, so their product divides $\Lprefix(x)$ as well.
\end{proof}

\noindent Formalised as the
\lword{Erdos269/ThreePrimeRunningLcm.lean}{124}{smoothPrefixLcm_eq_threePrimeHeight}{running-lcm
identity}, from the
\lword{Erdos269/ThreePrimeRunningLcm.lean}{60}{smooth3Val_dvd_threePrimeHeight_of_mem}{pointwise
divisibility}, the
\lword{Erdos269/ThreePrimeRunningLcm.lean}{78}{smoothPrefixLcm_dvd_threePrimeHeight}{divisibility
into the height} and the three pure-power memberships, namely the
\lword{Erdos269/ThreePrimeRunningLcm.lean}{85}{pureFirst_mem_smoothPrefixExponents}{first},
\lword{Erdos269/ThreePrimeRunningLcm.lean}{98}{pureSecond_mem_smoothPrefixExponents}{second}
and
\lword{Erdos269/ThreePrimeRunningLcm.lean}{111}{pureThird_mem_smoothPrefixExponents}{third}.
The unrestricted analogue is classical: iterating the prime-exponent maximum
rule gives $\lcm(1,\ldots,N)=\prod_{t\le N}t^{\lfloor\log_t N\rfloor}$ over the
primes $t\le N$~\cite[Ex.~1.21(a), p.~22]{apostol1976}, equivalently
$\log\lcm(1,\ldots,N)=\psi(N)$~\cite[\S4.2, p.~75]{apostol1976}, and
Montgomery and Vaughan state the identity directly in Exercise~6.2.7
\cite[p.~183]{montgomeryvaughan2007}.  Fan states the $P$-restricted form for
every $|P|$~\cite{fan2026comment}.  The pairwise distinctness hypothesis is
used exactly once, in the coprimality step.

At $(p,q,r)=(2,3,5)$ the first ten values are
\[
\begin{array}{c|cccccccccc}
x&1&2&3&4&5&6&7&8&9&10\\ \hline
\Lprefix(x)&1&2&6&12&60&60&60&120&360&360
\end{array}
\]
So $\Lprefix(6)=4\cdot3\cdot5=60$, which exceeds $6$: the running value at a
smooth cutoff already contains powers of the other two primes that the cutoff
itself does not.  That small fact is the mechanism behind
Section~\ref{long269:sec:rank}.  Also $\hgt(x)\le x^{3}$, since each of the three
factors is a power of its base not exceeding $x$; this is the
\lword{Erdos269/ThreePrimeRunningLcm.lean}{431}{threePrimeHeight_le_cube}{cubic
majorant}, and the exponent is the number of generating primes.

Say that $x$ and $y$ lie in the same \emph{logarithmic cell} when
$\lfloor\log_b x\rfloor=\lfloor\log_b y\rfloor$ for each of $b=p,q,r$, the
\lword{Erdos269/ThreePrimeRunningLcm.lean}{164}{SameThreePrimeLogCell}{cell
relation}.  By Theorem~\ref{long269:res:lcm} the running value depends on $x$ only
through the three integer logarithms, so it is constant on cells and moves only
where one logarithm moves.

\begin{proposition}[constancy and jump ratios]\label{long269:res:cell}
If $x,y\ge1$ lie in the same logarithmic cell then $\Lprefix(x)=\Lprefix(y)$,
and the same holds for the kernel at two smooth points of one cell.  If
$\lfloor\log_p y\rfloor=\lfloor\log_p x\rfloor+1$ while the other two
logarithms agree, then $\Lprefix(y)=p\,\Lprefix(x)$, and similarly with $q$ or
$r$ in place of $p$.
\end{proposition}

\begin{proof}
Both parts are immediate from Theorem~\ref{long269:res:lcm}: the height depends on $x$
only through the three integer logarithms, and advancing one of them multiplies
exactly one factor by its base.
\end{proof}

\noindent The first part is the
\lword{Erdos269/ThreePrimeRunningLcm.lean}{180}{smoothPrefixLcm_eq_of_sameLogCell}{cell
constancy of the running value}, over the
\lword{Erdos269/ThreePrimeRunningLcm.lean}{170}{threePrimeHeight_eq_of_sameLogCell}{cell
constancy of the height} and the
\lword{Erdos269/ThreePrimeRunningLcm.lean}{192}{threePrimeKernelQ_eq_of_sameLogCell}{cell
constancy of the kernel}.  The jump statement is
\label{long269:res:jump}the
\lword{Erdos269/ThreePrimeRunningLcm.lean}{327}{smoothPrefixLcm_firstLogStep}{first},
\lword{Erdos269/ThreePrimeRunningLcm.lean}{340}{smoothPrefixLcm_secondLogStep}{second}
and
\lword{Erdos269/ThreePrimeRunningLcm.lean}{353}{smoothPrefixLcm_thirdLogStep}{third}
coordinate steps, over the corresponding height steps, which need no primality:
the
\lword{Erdos269/ThreePrimeRunningLcm.lean}{295}{threePrimeHeight_firstLogStep}{first},
\lword{Erdos269/ThreePrimeRunningLcm.lean}{306}{threePrimeHeight_secondLogStep}{second}
and
\lword{Erdos269/ThreePrimeRunningLcm.lean}{317}{threePrimeHeight_thirdLogStep}{third}.

\begin{proposition}[jump count]\label{long269:res:count}
Let $n\ge0$.  The set of the first $n$ positive powers of $p$, of $q$ and of
$r$ has exactly $3n$ elements, and adjoining the common origin $1$ gives
exactly $3n+1$.
\end{proposition}

\begin{proof}
Within one channel the powers $b,b^{2},\ldots,b^{n}$ are distinct because
$b\ge2$.  Across two channels a common value would be a positive power of two
distinct primes, which unique factorisation forbids.  Finally $1$ is not a
positive power of any prime.
\end{proof}

\noindent Formalised as the
\lword{Erdos269/ThreePrimeRunningLcm.lean}{250}{threePrimePositiveJumpSet_card}{positive
jump count} and the
\lword{Erdos269/ThreePrimeRunningLcm.lean}{279}{threePrimeJumpSetWithOrigin_card}{jump
count with the origin}, over the
\lword{Erdos269/ThreePrimeRunningLcm.lean}{209}{positivePrimePowers_card}{channel
cardinality}, the
\lword{Erdos269/ThreePrimeRunningLcm.lean}{230}{positivePrimePowers_disjoint}{channel
disjointness} and the
\lword{Erdos269/ThreePrimeRunningLcm.lean}{220}{one_not_mem_positivePrimePowers}{exclusion
of the origin}; the channels are the
\lword{Erdos269/ThreePrimeRunningLcm.lean}{205}{positivePrimePowers}{positive
power sets}, whose union is the
\lword{Erdos269/ThreePrimeRunningLcm.lean}{244}{threePrimePositiveJumpSet}{positive
jump set} and, with the origin adjoined, the
\lword{Erdos269/ThreePrimeRunningLcm.lean}{274}{threePrimeJumpSetWithOrigin}{full
jump set}.  Two channels never meet, so at each positive pure power exactly
one logarithm advances and the running value is multiplied there by the prime
of that channel.  Reading those multipliers in increasing order of the pure
powers gives the \emph{jump word} of $\Lprefix$.  At $\{2,3,5\}$ the pure
powers in increasing order are $2,3,4,5,8,9,16,25,27,32,\ldots$, so the jump
word begins
\[
 2,\;3,2,\;5,2,\;3,2,\;5,3,2,\;\ldots
\]
grouped so that each group ends at a power of two, which is the grouping used
in Section~\ref{long269:sec:blocks}.

Two finite statements are recorded here because they are the exact shape of the
counting used later.  Write $\mathcal B(h_p,h_q,h_r)$ for the box of exponent
triples with $i\le h_p$, $j\le h_q$, $k\le h_r$, the
\lword{Erdos269/ThreePrimeRunningLcm.lean}{369}{smoothExponentBox}{exponent
box}, and $F(H)$ for the set of its points whose height equals $H$, the
\lword{Erdos269/ThreePrimeRunningLcm.lean}{378}{smoothHeightFiber}{height
fibre}.

\begin{proposition}[finite normal form]\label{long269:res:fibre}
For every box $\mathcal B$,
$\sum_{(i,j,k)\in\mathcal B}\Kernel(i,j,k)=\sum_{H}\#F(H)/H$, the outer sum
ranging over the heights attained on $\mathcal B$.
\end{proposition}

\begin{proof}
Partition $\mathcal B$ into the fibres of the height map.  On $F(H)$ every
summand is $1/H$, so the fibre contributes $\#F(H)/H$.
\end{proof}

\noindent Formalised as the
\lword{Erdos269/ThreePrimeRunningLcm.lean}{407}{finiteSmoothKernelSum_groupedByHeight}{height-fibre
normal form}, over the
\lword{Erdos269/ThreePrimeRunningLcm.lean}{385}{smoothHeightFiber_kernel_sum}{fibre
sum} and the
\lword{Erdos269/ThreePrimeRunningLcm.lean}{374}{smoothPointHeight}{point
height}.  At $(2,3,5)$ on $\mathcal B(1,1,1)$ the eight smooth values
$1,2,3,5,6,10,15,30$ carry heights $1,2,6,60,60,360,360,10800$, so two
coefficients equal $2$ and the identity reads $18421/10800$ on both sides.  The
identity is between two finite sums over a box, so it is not a statement about
$\Lprefix$ at a cutoff.  The same module records the one-step map
$\tau(b,d,s)=b(s-d)$, the
\lword{Erdos269/ThreePrimeRunningLcm.lean}{488}{tailStateStep}{variable-base
tail step}, with the rewriting
\lword{Erdos269/ThreePrimeRunningLcm.lean}{491}{tailStateStep_eq}{that names its
expanded form}; the orbits analysed in Section~\ref{long269:sec:blocks} are the actual
ones, and no orbit of this abstract map is analysed.

Now fix an interval $[\lo,\hi)$ and write $\mathcal S$ for the exponent triples
of $\mathcal B(h_p,h_q,h_r)$ whose smooth value lies in it, the
\lword{Erdos269/ThreePrimeRunningLcm.lean}{501}{smoothExponentShell}{smooth
exponent shell}.

\begin{lemma}[uniqueness in a short interval]\label{long269:res:short}
Let $b\ge1$ and $\hi\le b\,\lo$.  If $b^{a}w$ and $b^{a'}w$ both lie in
$[\lo,\hi)$ then $a=a'$.
\end{lemma}

\begin{proof}
If $a<a'$ then $\hi\le b\,\lo\le b^{a+1}w\le b^{a'}w<\hi$, which is impossible;
the case $a>a'$ is symmetric.
\end{proof}

\begin{proposition}[projection and the quadratic shell bound]\label{long269:res:drop}
\label{long269:res:shell}
If $\hi\le r\,\lo$ then $\#\mathcal S\le(h_p+1)(h_q+1)$, and if
$\hi\le p\,\lo$ then $\#\mathcal S\le(h_q+1)(h_r+1)$.  If moreover
$\hi\le r\,\lo$ and $h_p\le h_q\le h_r$ with $h_p+h_q+h_r=j$, then
$9\,\#\mathcal S\le(j+3)^{2}$.
\end{proposition}

\begin{proof}
Suppose $\hi\le r\,\lo$.  If two triples of $\mathcal S$ agree in their first
two coordinates, Lemma~\ref{long269:res:short} with $b=r$ and $w=p^{i}q^{j}$ forces
their third coordinates to agree, so the projection forgetting the third
coordinate is injective on $\mathcal S$ and its image lies in a rectangle with
$(h_p+1)(h_q+1)$ points.  The other case is the same with the first coordinate
projected away.  Under the sorting hypothesis the two surviving coordinates are
the two smallest, so it suffices that $a\le b\le c$ with $a+b+c=j$ gives
$9(a+1)(b+1)\le(j+3)^{2}$.  From $a\le b\le c$ we get $a+2b\le j$, so it is
enough that $9(a+1)(b+1)\le(a+2b+3)^{2}$; writing $b=a+d$ with $d\ge0$, the
difference of the two sides is $d(3a+4d+3)\ge0$.
\end{proof}

\noindent Formalised as the
\lword{Erdos269/ThreePrimeRunningLcm.lean}{586}{smoothExponentShell_card_le_dropThird}{third-coordinate
projection}, the
\lword{Erdos269/ThreePrimeRunningLcm.lean}{544}{smoothExponentShell_card_le_dropFirst}{first-coordinate
projection}, the
\lword{Erdos269/ThreePrimeRunningLcm.lean}{647}{smoothExponentShell_card_quadratic}{quadratic
shell bound} and the
\lword{Erdos269/ThreePrimeRunningLcm.lean}{630}{sorted_pair_quadratic}{sorted
quadratic estimate}, over the
\lword{Erdos269/ThreePrimeRunningLcm.lean}{510}{exponent_unique_in_short_interval}{short-interval
uniqueness}.  The constant $9$ is the square of the number of generating
primes, and the last inequality of the proof is an equality when
$h_p=h_q=h_r$, so no larger constant survives that step.  The estimate uses no
analytic input on the distribution of smooth numbers.

\section{One Hecke--Mahler value controls both two-prime sums}
\label{long269:sec:two-prime}

Temporarily let $P=\{p,q\}$ with $p<q$, and write
$L_{p,q}(t)=p^{\lfloor\log_p t\rfloor}q^{\lfloor\log_q t\rfloor}$, which is the
running least common multiple of the $\{p,q\}$-smooth numbers up to $t$ by the
argument of Theorem~\ref{long269:res:lcm} with one coordinate omitted.  The
de-duplicated sum retains the initial value $1$ and one reciprocal for every
later distinct running value, so
\[
 \mathcal D_{\{p,q\}}
 =1+\sum_{t\in\{p,p^2,\ldots\}\cup\{q,q^2,\ldots\}}\frac1{L_{p,q}(t)},
 \qquad
 \mathcal R_{\{p,q\}}=\sum_{i,j\ge0}\frac1{L_{p,q}(p^iq^j)} .
\]

\begin{proof}[Proof of Theorem~\ref{long269:res:lead-two-prime}]
Set
\[
 \theta=\frac{\log p}{\log q},\qquad x=\frac1p,\qquad y=\frac1q,\qquad
 m_n=\lfloor n\theta\rfloor,\qquad \delta_n=m_{n+1}-m_n .
\]
Here $0<\theta<1$, and $\theta$ is irrational, since a rational value would
give $p^{b}=q^{a}$ for positive integers $a,b$.  Consequently
$\delta_n\in\{0,1\}$.  Put
\[
 A=\sum_{n\ge0}x^ny^{m_n},
 \qquad
 B_\ast=\sum_{n\ge0}\delta_nx^ny^{m_n+1} .
\]
The initial value and the $p$-channel contribute $A$, since
$L_{p,q}(p^n)=p^nq^{m_n}$.  A power of $q$ lies strictly between $p^n$ and
$p^{n+1}$ exactly when $\delta_n=1$, it is then $q^{m_n+1}$, and its post-jump
reciprocal is $x^ny^{m_n+1}$; so the $q$-channel contributes $B_\ast$ and
$\mathcal D_{\{p,q\}}=A+B_\ast$.  All these series converge absolutely.

Since $y^{m_{n+1}}-y^{m_n}=\delta_ny^{m_n}(y-1)$, an index shift gives
$A-1-xA=x(y-1)B_\ast/y$, so $B_\ast=\bigl((p-1)A-p\bigr)/(1-q)$ and
\begin{equation}\label{long269:eq:two-prime-affine}
 \mathcal D_{\{p,q\}}=\frac{(q-p)A+p}{q-1} .
\end{equation}
At a smooth point, $\log_p(p^iq^j)=i+j/\theta$ and
$\log_q(p^iq^j)=j+i\theta$, so
$L_{p,q}(p^iq^j)=p^{\,i+\lfloor j/\theta\rfloor}q^{\,j+m_i}$ and absolute
convergence permits the factorisation
\[
 \mathcal R_{\{p,q\}}=AC,\qquad C=\sum_{j\ge0}y^jx^{\lfloor j/\theta\rfloor}.
\]
For $j\ge1$ put $n=\lfloor j/\theta\rfloor$.  Then $n>j/\theta-1$ gives
$n\theta>j-\theta>j-1$, and $n+1>j/\theta$ gives $(n+1)\theta>j$, while
$n\theta\le j$ and $(n+1)\theta<j+\theta<j+1$; hence $m_n=j-1$, $m_{n+1}=j$ and
$\delta_n=1$.  Conversely every $n$ with $\delta_n=1$ arises from the unique
$j=m_{n+1}$.  Therefore $C=1+B_\ast$ and
\begin{equation}\label{long269:eq:two-prime-quadratic}
 \mathcal R_{\{p,q\}}
 =\frac{(p+q-1)A-(p-1)A^{2}}{q-1} .
\end{equation}

It remains to prove that $A$ is transcendental.  For the Hecke--Mahler series
\[
 F_\theta(x,y)=\sum_{n\ge1}\sum_{k=1}^{\lfloor n\theta\rfloor}x^ny^k
\]
a finite geometric sum gives
$\bigl((1-y)/y\bigr)F_\theta(x,y)=x/(1-x)-(A-1)$, that is
\begin{equation}\label{long269:eq:hecke-mahler-boundary}
 A=\frac1{1-x}-\frac{1-y}{y}F_\theta(x,y).
\end{equation}
Bugeaud and Laurent's Theorem~1.1 states, in particular, that
$F_\theta(\beta,\alpha)$ is transcendental when $\theta\in(0,1)$ is irrational,
$\alpha$ and $\beta$ are nonzero algebraic numbers, $|\beta|<1$ and
$|\beta\alpha^\theta|<1$ \cite[p.~61, Theorem~1.1]{bugeaudlaurent2023}; the
$\rho=0$ case used here goes back to Loxton and van der Poorten
\cite[p.~40, Theorem~8]{loxtonvdp1977}.  Take $(\beta,\alpha)=(x,y)$: then
$|\beta|=1/p<1$ and
\[
 |xy^{\theta}|=\frac1p\left(\frac1q\right)^{\log p/\log q}=\frac1{p^{2}}<1 .
\]
So $F_\theta(x,y)$ is transcendental, and \eqref{long269:eq:hecke-mahler-boundary}
makes $A$ transcendental.  The coefficient of $A$
in~\eqref{long269:eq:two-prime-affine} is $(q-p)/(q-1)\ne0$ and the coefficient of
$A^{2}$ in~\eqref{long269:eq:two-prime-quadratic} is $-(p-1)/(q-1)\ne0$, both rational.
If either value were algebraic, its identity would exhibit $A$ as a root of a
nonzero polynomial over the algebraic numbers.
\end{proof}

\noindent The two identities are separately labelled
\label{long269:res:two-prime-transcendence}\label{long269:res:two-prime-repeated-transcendence}
because they carry different weight.  A product of two transcendental numbers
need not be transcendental, so the factorisation $\mathcal R=AC$ alone proves
nothing; the quadratic identity~\eqref{long269:eq:two-prime-quadratic} in the single
value $A$ is what settles the repeated sum.  A third prime replaces the single
Beatty boundary by a genuinely two-dimensional ordering problem, and the next
section makes that failure exact.

\section{The rank phase transition}\label{long269:sec:rank}

One might hope to write the kernel as $f(i)g(j)h(k)$ and reduce the problem to
one-dimensional criteria.  At two generators that hope is exactly correct, and
at three it fails at every finite order.

\begin{proposition}[two generators separate]\label{long269:res:two-prime-rank}
For $p,q>1$ and all $i,j$, the two-prime kernel
$\Kernel_2(i,j)=1/L_{p,q}(p^iq^j)$ is the outer product
\[
 \Kernel_2(i,j)
 =\bigl(p^{i}q^{\lfloor\log_q p^{i}\rfloor}\bigr)^{-1}
  \bigl(p^{\lfloor\log_p q^{j}\rfloor}q^{j}\bigr)^{-1},
\]
so every two-by-two minor of $\Kernel_2$ vanishes.
\end{proposition}

\begin{proof}
By the two-variable form of Theorem~\ref{long269:res:lcm} the exponents of $p$ and of
$q$ in $L_{p,q}(p^iq^j)$ are $i+\lfloor\log_p q^{j}\rfloor$ and
$j+\lfloor\log_q p^{i}\rfloor$, each depending on one index only.  A matrix
whose entries are a product of a row function and a column function has
vanishing two-by-two minors.
\end{proof}

\noindent Checked as the
\lword{Erdos269/KernelCarryRank.lean}{278}{twoPrimeKernelQ_eq_outer_product}{outer-product
identity} and the
\lword{Erdos269/KernelCarryRank.lean}{298}{twoPrimeKernelQ_minor_two_eq_zero}{vanishing
of every two-by-two minor}.  This is the linear-algebraic shadow of
Section~\ref{long269:sec:two-prime}: the two-prime value factors because the kernel
does.

At three generators the smallest rectangle already fails to factor.  A
factorisation $f(i)g(j)h(k)$ would force
$\Kernel(0,0,0)\Kernel(1,1,0)=\Kernel(1,0,0)\Kernel(0,1,0)$.

\begin{proposition}[non-separability at $\{2,3,5\}$]\label{long269:res:rank}
With $(p,q,r)=(2,3,5)$,
\[
 \det\begin{pmatrix}
 \Kernel(0,0,0)&\Kernel(0,1,0)\\
 \Kernel(1,0,0)&\Kernel(1,1,0)
 \end{pmatrix}
 =\det\begin{pmatrix}1&1/6\\1/2&1/60\end{pmatrix}
 =-\frac1{15}\ne0 .
\]
\end{proposition}

\begin{proof}
The four values are computed from $\hgt(1)=1$, $\hgt(2)=2$,
$\hgt(3)=2\cdot3=6$ and $\hgt(6)=4\cdot3\cdot5=60$, so the determinant is
$1/60-1/12=-1/15$.
\end{proof}

\noindent Checked as the
\lword{Erdos269/ThreePrimeRunningLcm.lean}{827}{kernel_235_minor_eq_neg_one_fifteen}{rank-two
certificate} and the
\lword{Erdos269/ThreePrimeRunningLcm.lean}{480}{kernel_235_not_rankOne}{non-separation
witness}, over the four exact values, the
\lword{Erdos269/ThreePrimeRunningLcm.lean}{444}{kernel_235_origin}{origin},
\lword{Erdos269/ThreePrimeRunningLcm.lean}{451}{kernel_235_two}{value at two},
\lword{Erdos269/ThreePrimeRunningLcm.lean}{459}{kernel_235_three}{value at
three} and
\lword{Erdos269/ThreePrimeRunningLcm.lean}{468}{kernel_235_six}{value at six};
the third-order witness $1/81000$ is
\lword{Erdos269/KernelCarryRank.lean}{570}{kernel_235_minor3_eq_one_over_81000}{checked}
beside the
\lword{Erdos269/KernelCarryRank.lean}{560}{kernel_235_minor2_eq_neg_one_fifteen}{second-order
one}.  The height at $6$ is $60$ because the maximal pure powers below $6$ are
$4$, $3$ and $5$, so the running value at a smooth cutoff sees powers of the
other primes that the cutoff itself omits.

\begin{theorem}[arbitrary-order non-separability]\label{long269:res:infinite-rank}
\label{long269:res:lead-infinite-rank}
Let $p,q,r$ be primes with $p\ne q$, $p\ne r$ and $q\ne r$.  For every
$n\ge1$ there are injective maps $I,J:\{0,\ldots,n-1\}\to\N$ such that, for
every $k\ge0$,
\[
 \det\bigl(\Kernel(I(a),J(b),k)\bigr)_{0\le a,b<n}\ne0 .
\]
Consequently, for no finite $d$ do there exist rational-valued functions
$f_\ell:\N\to\Q$ and $G_\ell:\N^{2}\to\Q$, $0\le\ell<d$, satisfying
$\Kernel(i,j,k)=\sum_{\ell<d}f_\ell(i)G_\ell(j,k)$ for all $i,j,k$.
\end{theorem}

\begin{proof}
Put $\alpha=\log_r p$, $\beta=\log_r q$, $x_i=\{i\alpha\}$ and
$y_j=\{j\beta\}$.  The three height exponents are
\[
 \begin{split}
 v_p(\hgt(p^iq^jr^k))&=i+\lfloor j\log_pq+k\log_pr\rfloor,\\
 v_q(\hgt(p^iq^jr^k))&=j+\lfloor i\log_qp+k\log_qr\rfloor,\\
 v_r(\hgt(p^iq^jr^k))&=k+\lfloor i\alpha\rfloor+\lfloor j\beta\rfloor+
   \lfloor x_i+y_j\rfloor .
 \end{split}
\]
Every term except the last floor depends on $i$ and $k$ alone or on $j$ and $k$
alone, so with positive rational $R_i(k)$ and $C_j(k)$,
\begin{equation}\label{long269:eq:carry-factorisation}
 \Kernel(i,j,k)=R_i(k)\,C_j(k)\,t^{\lfloor x_i+y_j\rfloor},
 \qquad t=r^{-1} .
\end{equation}
The remaining matrix is independent of $k$, and $x_i+y_j\in[0,2)$, so its
entries are $1$ and $t$.

Each of $\alpha$ and $\beta$ is irrational, since a rational value would give
an equality of positive powers of distinct primes, so each fractional-part
orbit is dense in $(0,1)$ and each is injective.  Choose indices with
$0<x_{I(0)}<\cdots<x_{I(n-1)}<1$, then $y_{J(0)}$ in $(1-x_{I(0)},1)$ and, for
$b>0$, $y_{J(b)}$ in $(1-x_{I(b)},1-x_{I(b-1)})$.  These intervals are
disjoint, so both index maps are injective, and $x_{I(a)}+y_{J(b)}\ge1$ exactly
when $b\le a$.  The remaining matrix is therefore
\[
 C_n(t)=\begin{pmatrix}
 t&1&1&\cdots&1\\
 t&t&1&\cdots&1\\
 \vdots&\vdots&\ddots&\ddots&\vdots\\
 t&t&\cdots&t&1\\
 t&t&\cdots&t&t
 \end{pmatrix},
 \qquad
 \det C_n(t)=t(t-1)^{n-1}\ne0 ,
\]
the determinant following by subtracting from each row its predecessor, working
upwards from the last.  Equation~\eqref{long269:eq:carry-factorisation} multiplies this
determinant by nonzero row and column factors, for every $k$, which is the
asserted uniformity.  Finally, a representation with $d$ summands would factor
every $(d+1)\times(d+1)$ restriction through a $d$-dimensional space and force
its determinant to vanish.
\end{proof}

\noindent The factorisation~\eqref{long269:eq:carry-factorisation} is the
\lword{Erdos269/KernelCarryRank.lean}{259}{threePrimeKernelQ_factorisation}{checked
kernel factorisation}, over the
\lword{Erdos269/KernelCarryRank.lean}{218}{threePrimeHeight_factorisation}{height
factorisation}.  The uniform minor family is
\lword{Erdos269/KernelCarryRank.lean}{376}{exists_uniform_nonsingular_threePrimeKernel_minor_of_prime}{checked},
the exclusion of every finite separation is
\lword{Erdos269/KernelCarryRank.lean}{388}{not_finite_separable_threePrimeKernel}{checked},
and the index selection is discharged by the problem-neutral staircase engine
\lword{Shared/IrrationalRotationStaircase.lean}{271}{exists_staircase_indices}{exists\_staircase\_indices},
which realises the full $n\times n$ pattern from two rotations without integer
returns and depends on Mathlib only; the absence of integer returns is the
\lword{Erdos269/KernelCarryRank.lean}{154}{noIntegerOrbit_logb_of_prime}{checked
irrationality of the two logarithm ratios}.  The Lean forms assume that $p$,
$q$ and $r$ are prime with $p\ne r$ and $q\ne r$, and they do not use $p\ne q$.
The proof uses the two one-dimensional density statements separately and needs
no joint equidistribution hypothesis.

Index selection is essential.  The leading minors of the
$\{2,3,5\}$ kernel are not a witness: row three is $1/120$ times row zero for
$j\le3$
(\lword{Erdos269/KernelCarryRank.lean}{583}{kernel_235_row_three_eq_smul_row_zero}{checked}),
and the proportionality fails at $j=4$
(\lword{Erdos269/KernelCarryRank.lean}{592}{kernel_235_row_three_ne_smul_row_zero_at_four}{checked}).
A proof that read off leading minors alone would be false.

\paragraph{Attribution and reach.}
Fan recorded the qualitative obstruction for $|P|\ge3$ in his comment of
26 June 2026~\cite{fan2026comment}, observing that the two-prime argument does
not seem to generalise immediately.  The quantitative form above, its
uniformity in the third coordinate, and its formalisation are proved here.
Theorem~\ref{long269:res:infinite-rank} excludes exact separations of the displayed
finite-sum form; it is not an independence, irrationality or transcendence
statement, and it settles no instance of Erd\H{o}s~\#269.

\subsection{A sharp uniform obstruction for the normalised carry matrix}

The determinant argument is exact and therefore fragile: an arbitrarily small
perturbation of the kernel can destroy an exact vanishing.  The next theorem
replaces exactness by a uniform distance, and the constant it produces is
sharp.  It concerns the normalised carry matrix
\begin{equation}\label{long269:eq:carry-matrix}
 C(i,j)=t^{\lfloor x_i+y_j\rfloor},
 \qquad
 x_i=\{i\log_r p\},\quad y_j=\{j\log_r q\},\quad t=r^{-1},
\end{equation}
which is the factor left in~\eqref{long269:eq:carry-factorisation} after the row and
column factors are divided out.  Say that a real matrix $A$ on $\N\times\N$ has
\emph{finite separated rank} when all of its columns lie in one
finite-dimensional space of real sequences, equivalently when
$A(i,j)=\sum_{\ell<d}f_\ell(i)g_\ell(j)$ for some finite $d$ and some
sequences $f_\ell,g_\ell$, with no continuity or boundedness assumed.

\begin{theorem}[sharp uniform separated approximation]\label{long269:res:uniform-rank}
Let $p,q,r$ be pairwise distinct primes and let $C$ be as
in~\eqref{long269:eq:carry-matrix}.  Then
\[
 \inf_{A}\ \sup_{i,j\ge0}\ |C(i,j)-A(i,j)|=\frac{1-t}{2}=\frac{r-1}{2r},
\]
the infimum being over all matrices $A$ of finite separated rank, and it is
attained by the constant matrix of value $(1+t)/2$.
\end{theorem}

\begin{proof}
Every entry of $C$ lies in $\{1,t\}$, and $C(i,j)=t$ exactly when
$x_i+y_j\ge1$.  Fix $j\ne k$.  The numbers $y_j$ are pairwise distinct, since
$\log_r q$ is irrational, so we may assume $y_j<y_k$, and then
$0\le1-y_k<1-y_j\le1$.  Density of the orbit $(x_i)$ in $(0,1)$ supplies an
index $i$ with $1-y_k<x_i<1-y_j$, and at that row the two columns carry the
entries $t$ and $1$.  Hence any two distinct columns of $C$ are at sup-distance
exactly $1-t$.

Let $A$ have finite separated rank and put
$\varepsilon=\sup_{i,j}|C(i,j)-A(i,j)|$.  Suppose $\varepsilon<(1-t)/2$.  Each
column $A_j$ then satisfies $\|A_j\|_\infty\le1+\varepsilon$, so all columns of
$A$ lie in $W=V\cap\ell^{\infty}$, where $V$ is the finite-dimensional space
containing the columns of $A$; the space $W$ is a subspace of $V$, hence
finite-dimensional, and $\|\cdot\|_\infty$ is a genuine norm on it.  By the
triangle inequality, distinct columns satisfy
\[
 \|A_j-A_k\|_\infty\ \ge\ \|C_j-C_k\|_\infty-2\varepsilon
 \ =\ 1-t-2\varepsilon\ >\ 0 .
\]
The compactness input is the standard fact that closed bounded subsets of
a finite-dimensional real normed space are compact. The use of
$W=V\cap\ell^\infty$ is essential: the individual separated row factors
need not be bounded.
The columns $A_j$, $j\ge0$, are therefore infinitely many points of $W$, all of
norm at most $1+\varepsilon$ and pairwise separated by a fixed positive
distance.  A bounded subset of a finite-dimensional normed space is totally
bounded, so it contains no infinite uniformly separated family.  This
contradiction gives $\varepsilon\ge(1-t)/2$ for every $A$ of finite separated
rank.

For sharpness take $A(i,j)=(1+t)/2$, which has separated rank one; every entry
of $C$ is at distance exactly $(1-t)/2$ from it.
\end{proof}

The compactness step is standard, and the statement is recorded here for this
kernel without a claim of technique.  Three qualifications fix its
reach.  It is a statement about the normalised carry
matrix~\eqref{long269:eq:carry-matrix}: the original kernel carries positive row and
column factors that decay, so the conclusion transfers to a weighted uniform
norm relative to those factors and not to the unweighted sup norm of
$\Kernel$.  It bounds approximation of the whole infinite matrix, and on any
finite range a separated approximant of small rank exists.  And it says nothing
about approximating the scalar sum, so it yields no irrationality conclusion.
What it does add to Theorem~\ref{long269:res:infinite-rank} is robustness: no finite
separated model of the carry, however chosen, gets uniformly closer than half a
carry jump.  The threshold is exact in both directions, since rank one attains
it.

\section{Dyadic blocks and the literal infinite tail}\label{long269:sec:blocks}

For the rest of the note set $P=\{2,3,5\}$ and $S=\mathcal R_P$, and write
$\hgt_a=\hgt(2^{a})$ and $h_a=\hgt_a/2$, so that $h_0=1/2$ while $h_a$ is a
positive integer for $a\ge1$.  Define the dyadic shell mass, its tail and its
normalisation by
\begin{equation}\label{long269:eq:actual-tail}
 s_a=\sum_{\substack{i,j,k\ge0\\ 2^{a}\le2^{i}3^{j}5^{k}<2^{a+1}}}
       \frac1{\hgt(2^{i}3^{j}5^{k})},
 \qquad
 U_a=\sum_{j\ge a}s_j,
 \qquad
 X_a=h_aU_a .
\end{equation}

Compress the jump word of Section~\ref{long269:sec:lcm} into the blocks cut out by
consecutive powers of two: block $a$ starts just after $2^{a}$, contains every
pure $3$- or $5$-power strictly between $2^{a}$ and $2^{a+1}$, and ends with
the jump at $2^{a+1}$.  A channel cannot occur twice inside one block, because
two powers of the same base $b\ge2$ inside an interval of ratio $2\le b$ must
coincide; this is the
\lword{Erdos269/ThreePrimeRunningLcm.lean}{669}{dyadicInternalPower_exponent_unique}{checked
internal-power uniqueness lemma} for the
\lword{Erdos269/ThreePrimeRunningLcm.lean}{663}{DyadicInternalPower}{internal-power
predicate}.  Let $I_a$ list the internal jumps $(p,e)$
with $p\in\{3,5\}$ and $2^{a}<p^{e}<2^{a+1}$, in increasing order.

\begin{proposition}[the dyadic block alphabet]\label{long269:res:dyadic-alphabet}
For every $a$,
\begin{equation}\label{long269:eq:dyadic-alphabet}
 b_a=\frac{\hgt_{a+1}}{\hgt_a}=2\prod_{(p,e)\in I_a}p\in\{2,6,10,30\},
 \qquad\text{so}\qquad 2\le b_a\le30 .
\end{equation}
\end{proposition}

\begin{proof}
Internal-power uniqueness leaves two independent yes-or-no choices, one for the
$3$-channel and one for the $5$-channel.  Multiplying the terminal factor $2$
by the selected channel factors gives exactly the four displayed cases.
\end{proof}

\noindent The definition is the
\lword{Erdos269/ThreePrimeRunningLcm.lean}{751}{dyadicBlockBase235}{dyadic
block base}, and Lean checks both the
\lword{Erdos269/ThreePrimeRunningLcm.lean}{762}{dyadicBlockBase235_cases}{exact
four-case alphabet} and the
\lword{Erdos269/ThreePrimeRunningLcm.lean}{774}{dyadicBlockBase235_mem_interval}{bounded-radix
consequence}.  All four letters occur early:
\[
\begin{array}{c|c|c|c|c}
a & (2^{a},2^{a+1}) & \text{internal }3\text{-power} & \text{internal }5\text{-power} & b_a\\ \hline
0 & (1,2)   & \text{none} & \text{none} & 2\\
1 & (2,4)   & 3           & \text{none} & 6\\
2 & (4,8)   & \text{none} & 5           & 10\\
3 & (8,16)  & 9           & \text{none} & 6\\
4 & (16,32) & 27          & 25          & 30\\
5 & (32,64) & \text{none} & \text{none} & 2
\end{array}
\]

For $p\in P$ with $q,r$ the other two primes, put
\[
 A_p(e)=\#\{(i,j)\in\N^2:q^{i}r^{j}<p^{e}\},\qquad
 C_p(e)=\sum_{u=1}^{e}A_p(u),\qquad C_p(0)=0 ,
\]
and for $(p,e)\in I_a$ let $\sigma_a(p,e)$ be the product of the channels of
the later internal jumps of block $a$.  The \emph{ordered block digit} is
\begin{equation}\label{long269:eq:actual-digit}
 m_a=A_2(a+1)+\sum_{(p,e)\in I_a}(p-1)\,\sigma_a(p,e)\,
        \bigl(C_p(e)-C_2(a)\bigr)\quad(a\ge1),
 \qquad m_0=1 ,
\end{equation}
the integer implemented by the pinned checker and the
\lword{Erdos269/DyadicBlockThresholdPartition.lean}{150}{dyadicOrderedBlockDigit235}{checked
ordered digit}.  The first four pairs $(b_a,m_a)$ from $a=1$ are
$(6,4)$, $(10,7)$, $(6,7)$, $(30,65)$.

\begin{theorem}[the literal shell recurrence]\label{long269:res:actual-orbit}
The shell masses are summable and $S=\sum_{a\ge0}s_a$.  For every $a\ge0$,
\begin{equation}\label{long269:eq:shell-digit-identity}
 m_a=h_{a+1}s_a\in\N_{>0},
 \qquad
 X_{a+1}=b_aX_a-m_a,
 \qquad
 X_a=\sum_{j\ge a}\frac{m_j}{b_ab_{a+1}\cdots b_j} .
\end{equation}
\end{theorem}

\begin{proof}
Every exponent of $3$ or $5$ in the $a$th shell is at most $a$, and for each
such pair Lemma~\ref{long269:res:short} leaves at most one exponent of $2$ placing the
point in $[2^{a},2^{a+1})$.  Each height is at least $2^{a}$ and the shell
contains $2^{a}$, so $0<s_a\le(a+1)^{2}2^{-a}$.  The majorant is summable,
which justifies every tail splitting below, and unique prime factorisation
identifies $\sum_a s_a$ with the original repeated series.

Write the internal jumps as $t_\ell=p_\ell^{e_\ell}$, $1\le\ell\le v$, and set
$t_0=2^{a}$, $t_{v+1}=2^{a+1}$.  Let $N(t)$ count smooth positive integers
strictly below $t$ and put $n_\ell=N(t_\ell)$.  On $[t_\ell,t_{\ell+1})$ the
height is $\hgt_a\prod_{u\le\ell}p_u$, and the terminal jump has factor two, so
\[
 h_{a+1}s_a
 =\sum_{\ell=0}^{v}(n_{\ell+1}-n_\ell)\prod_{u>\ell}p_u
 =n_{v+1}-n_0+\sum_{\ell=1}^{v}(p_\ell-1)
        \Bigl(\prod_{u>\ell}p_u\Bigr)(n_\ell-n_0),
\]
the second equality being finite summation by parts.  Counting by the exponent
of $p$ gives $N(p^{e})=C_p(e)$, so the right-hand side is exactly
\eqref{long269:eq:actual-digit}; at $a=0$ the shell is the single point $1$, giving
$m_0=1$.  Positivity follows from $s_a>0$.  Splitting $U_a=s_a+U_{a+1}$ and
multiplying by $h_a$ gives the recurrence, and
$b_a\cdots b_j=\hgt_{j+1}/\hgt_a$ with $m_j=h_{j+1}s_j$ gives
$m_j/(b_a\cdots b_j)=h_as_j$, whose summation is the last identity.
\end{proof}

\noindent Summability is
\lword{Erdos269/DyadicShellSummability.lean}{142}{summable_dyadicShellMassR235}{checked},
and so is the
\lword{Erdos269/DyadicShellSummability.lean}{177}{dyadicNormalizedShellTsumTailR235_succ}{recurrence
for the genuine infinite tail}, over the
\lword{Erdos269/DyadicOrderedTailRecurrence.lean}{110}{dyadicNormalizedTailStateR235_succ}{normalised
state step}.

\begin{proposition}[integral state or cofinal separation]\label{long269:res:actual-dichotomy}
Either $X_a\in\Z$ for some $a\ge0$, or for every $a_0$ there is $a\ge a_0$ with
$|X_a-z|\ge1/31$ for every $z\in\Z$.
\end{proposition}

\begin{proof}
Suppose cofinal separation fails, so for some $A$ and every $a\ge A$ there is
an integer $z_a$ with $e_a=X_a-z_a$ and $|e_a|<1/31$.  The recurrence gives
$z_{a+1}-b_az_a+m_a=b_ae_a-e_{a+1}$.  The left side is an integer and the right
side has absolute value below $(30+1)/31=1$, so both vanish and
$e_{a+1}=b_ae_a$.  Hence $|e_{A+k}|\ge2^{k}|e_A|$ for every $k$ while
$|e_{A+k}|<1/31$, which forces $e_A=0$.
\end{proof}

\noindent This is
\lword{Erdos269/DyadicShellSummability.lean}{188}{dyadicShellTsumTail_integer_or_cofinal_far}{checked}
for the genuine infinite tail, over the abstract
\lword{Erdos269/BoundedRadixTailEscape.lean}{89}{boundedRadix_zero_or_cofinal_far}{integral-state
or cofinal-distance alternative}.  It excludes neither branch, and it says
nothing about a rational value whose reduced denominator has a factor coprime
to $30$: that case is the subject of Section~\ref{long269:sec:actual-orbit}.

\section{A quadratic tail bound and denominator cancellation}
\label{long269:sec:actual-orbit}

Put
\[
 n_a=a+\lfloor\log_3(2^{a})\rfloor+\lfloor\log_5(2^{a})\rfloor,
 \qquad
 Q(n)=\frac{n^{2}+8n+18}{9},
\]
so $n_a$ is the sum of the three height exponents at $2^{a}$.

\begin{theorem}[quadratic bound for the actual tail]\label{long269:res:actual-tail-bound}
For every $a\ge0$, $0<X_a\le Q(n_a)$.
\end{theorem}

\begin{proof}
Partition the smooth integers $x\ge2^{a}$ by their height vector
$(A,B,C)=(\lfloor\log_2x\rfloor,\lfloor\log_3x\rfloor,\lfloor\log_5x\rfloor)$.
The cell of a vector is the interval $[\lo,\hi)$ with
$\lo=\max(2^{A},3^{B},5^{C})$ and $\hi=\min(2^{A+1},3^{B+1},5^{C+1})$, so
$\hi\le2^{A+1}\le2\lo$ and, by Lemma~\ref{long269:res:short}, fixing the exponents of
$3$ and $5$ leaves at most one exponent of $2$.  Since $A\ge B\ge C$, the cell
contains at most
\[
 (B+1)(C+1)\le\frac{(A+B+C+3)^{2}}{9}
\]
smooth points: writing $u=B+1\ge v=C+1$ and using $A+1\ge u$, the difference
$(2u+v)^{2}-9uv=(u-v)(4u-v)$ is nonnegative.  Unique factorisation identifies
these exponent triples with distinct smooth integers.

As the cutoff increases all three height exponents are nondecreasing, so two
nonempty cells with the same exponent sum are the same cell, and there is at
most one nonempty cell of each rank.  Every cell above $2^{a}$ has rank at
least $n_a$, and a cell of rank $n_a+k$ has height at least $\hgt_a2^{k}$,
since each of the $k$ extra prime factors is at least $2$.  Nonnegative
summation over ranks, allowing empty ones, gives
\[
 X_a\le\frac1{18}\sum_{k\ge0}\frac{(n_a+k+3)^{2}}{2^{k}}
      =\frac{n_a^{2}+8n_a+18}{9},
\]
the evaluation using the geometric moments $\sum2^{-k}=2$, $\sum k2^{-k}=2$
and $\sum k^{2}2^{-k}=6$.  Positivity follows from the shell at $2^{a}$.
\end{proof}

The rank here is the sum of the three fixed height exponents, not a
varying-smoothness density parameter. The estimate uses only unique
factorisation, the short-cell counting bound and geometric moments; it does
not invoke a Dickman--Hildebrand asymptotic or a Hecke--Mahler theorem.

This is the analytic content of the section: each additional height rank costs
a geometric factor while its multiplicity grows only quadratically.  The finite
projection and the sorted quadratic inequality are the checked ingredients of
Proposition~\ref{long269:res:drop}; grouping the actual infinite tail by height cells
and summing the majorant is the argument above.

The bound is used at the endpoint of a window, where the natural index is the
positive prime-power jump count strictly below the cutoff.  Put
\begin{equation}\label{long269:eq:endpoint-index}
 j_a=\#\{p^{e}<2^{a}:p\in\{2,3,5\},\ e\ge1\}=n_a-1 ,
\end{equation}
the equality holding because the powers of $2$ below $2^{a}$ number $a-1$ while
the powers of $3$ and of $5$ below $2^{a}$ number $\lfloor\log_3 2^{a}\rfloor$
and $\lfloor\log_5 2^{a}\rfloor$.  Substituting $n_a=j_a+1$ into $Q$ gives the
integer bound used throughout the rest of the note:
\begin{equation}\label{long269:eq:actual-bound}
 K(B,a)=\bigl\lfloor B\,Q(n_a)\bigr\rfloor
 =\left\lfloor\frac{B(j_a^{2}+10j_a+27)}9\right\rfloor .
\end{equation}
The cutoff in \eqref{long269:eq:endpoint-index} is $2^{a}$ and it is strict.

\begin{lemma}[finite denominator clearing]\label{long269:res:all-scale-lattice}
For all integers $0\le u\le b$ the window mass $h_b\sum_{a=u}^{b-1}s_a$ is a
natural number.  If $S=N/D$ with $N\in\Z$ and $D\in\N_{>0}$, then $DX_a\in\Z$
for every $a\ge1$, and two states $X_i$, $X_j$ with $1\le i<j$ differ by an
integer.
\end{lemma}

\begin{proof}
An empty window has mass zero.  Otherwise $b\ge1$, and every integer $x<2^{b}$
has $2$-height exponent at most $b-1$ while its other height exponents are at
most those at $2^{b}$, so $\hgt(x)\mid\hgt_b/2=h_b$ and every term of the
finite window clears at $h_b$.  Writing $v_a=h_a\sum_{u<a}s_u\in\N$ gives
$DX_a=h_aN-Dv_a\in\Z$.  Among $D+1$ of these integers two share a residue
modulo $D$, and the corresponding states differ by an integer.
\end{proof}

\noindent The strict upper endpoint is what permits the division by two.
Checked as the
\lword{Erdos269/RationalLatticeReduction.lean}{158}{heightNormalizer235_mul_windowMass_eq_int}{window
clearing identity}, the
\lword{Erdos269/RationalLatticeReduction.lean}{196}{qsmul_normalizedTailState_eq_int_of_value_eq_rat}{all-scale
lattice} and the
\lword{Erdos269/RationalLatticeReduction.lean}{296}{exists_normalizedTailState_collision_of_value_eq_rat}{collision},
over the
\lword{Erdos269/RationalLatticeReduction.lean}{51}{smoothHeight_mul_prime_dvd_boundaryHeight}{boundary
divisibility} and the
\lword{Erdos269/RationalLatticeReduction.lean}{101}{two_mul_heightNormalizer235}{half-height
identity}.

\begin{theorem}[positive reduced carries from rationality]
\label{long269:res:actual-cancellation}\label{long269:res:lead-carry-bridge}
\label{long269:res:actual-carry-bound}\label{long269:res:denominator-reduction}
Suppose $S=N/D$ with $N\in\Z$, $D\in\N_{>0}$, and write
\[
 D=2^{u}3^{v}5^{w}B,\qquad u,v,w\in\N,\quad B\in\N_{>0},\quad\gcd(B,30)=1,
 \qquad a_D=u+1+2v+3w .
\]
Then for every $a\ge a_D$ the number $z_a=BX_a$ is a positive integer and
\[
 z_{a+1}=b_az_a-Bm_a,\qquad 1\le z_a\le K(B,a)\le90B(a+1)^{2} .
\]
\end{theorem}

\begin{proof}
Let $M=2^{u}3^{v}5^{w}$.  For $a\ge a_D$ we have $2^{a}\ge2^{u+1}$,
$2^{a}\ge3^{v}$ and $2^{a}\ge5^{w}$, using $3<2^{2}$ and $5<2^{3}$; hence
$M\mid h_a$.  In the clearing identity $DX_a=h_aN-Dv_a$ of
Lemma~\ref{long269:res:all-scale-lattice} both terms on the right are divisible by $M$,
so dividing by $M$ shows that $BX_a$ is an integer.  Positivity and the
recurrence come from~\eqref{long269:eq:shell-digit-identity}, and the upper bound is
Theorem~\ref{long269:res:actual-tail-bound} with the floor taken, since $z_a$ is an
integer below $BQ(n_a)$.  Finally $n_a\le3a$, so
$Q(n_a)\le a^{2}+\tfrac83a+2\le90(a+1)^{2}$.
\end{proof}

\noindent The smooth part $M$ affects only the onset $a_D$; the surviving carry
bound depends on the coprime factor $B$.  The bridge is
\lword{Erdos269/RationalityCarryBridge.lean}{324}{exists_reducedCarry_of_value_eq_rat}{checked}
against the cruder width $90B(a+1)^{2}$, which is the form the formal consumer
uses; the particular onset $a_D=u+1+2v+3w$ and the sharper bound
$K(B,a)$ are the paper statement above.  The abstract cancellation is checked
separately: every fixed smooth factor divides the running height once the
cutoff reaches it
(\lword{Erdos269/RestrictedFloorSum.lean}{548}{smooth3Val_dvd_threePrimeHeight_of_le}{absorption}),
a shared factor cancels from the recurrence
(\lword{Erdos269/RestrictedFloorSum.lean}{581}{integralCarry_cancel_commonFactor}{cancellation}),
and positivity, the bound and the window identity descend through it
(\lword{Erdos269/RestrictedFloorSum.lean}{601}{reducedCarry_pos_le_of_commonFactor}{bound
transfer},
\lword{Erdos269/RestrictedFloorSum.lean}{612}{reducedIntegralCarry_window}{window
transfer}).

The integral branch has one further exact property, recorded because it
constrains any attempt to build a surviving seed by hand.

\begin{proposition}[upward closure and rigidity]\label{long269:res:pinning}
For every $a$, $X_a=(m_a+X_{a+1})/b_a>0$, and if $X_a\in\Z$ then $X_n\in\Z$ for
every $n\ge a$.  Moreover, fix $A$, a positive width function $w$ with
$w(A+k)/2^{k}\to0$, and a real orbit $(y_n)_{n\ge A}$ satisfying
$y_{n+1}=b_ny_n-m_n$.  If $y_n$ and $X_n$ both lie in
$(m_n/b_n,\;m_n/b_n+w(n)]$ for every $n\ge A$, then $y_A=X_A$.
\end{proposition}

\begin{proof}
The identity is the recurrence solved for $X_a$, and positivity holds because
every shell contains its dyadic left endpoint.  Integer coefficients give
upward closure.  For the last assertion, the difference of the two orbits is
multiplied by $b_n\ge2$ at each step while the common window bounds its
absolute value at time $A+k$ by $w(A+k)$, so $2^{k}|y_A-X_A|\le w(A+k)$ and the
stated decay forces equality.
\end{proof}

\section{Residue windows and the equivalence band}\label{long269:sec:escape}

The statements of this section are about integer sequences.  For a start
$\ell\ge1$ and a length $h\ge1$ define
\[
 W_{\ell,h}=\prod_{j=0}^{h-1}b_{\ell+j},
 \qquad
 F_{\ell,0}=0,
 \qquad
 F_{\ell,h+1}=b_{\ell+h}F_{\ell,h}+m_{\ell+h},
\]
the
\lword{Erdos269/RestrictedFloorSum.lean}{417}{windowBase}{window base} and the
\lword{Erdos269/RestrictedFloorSum.lean}{422}{windowForcing}{window forcing}.
Iterating the recurrence gives the division-free identity
\begin{equation}\label{long269:eq:window-identity}
 X_{\ell+h}=W_{\ell,h}X_\ell-F_{\ell,h},
 \qquad\text{and}\qquad
 z_{\ell+h}=W_{\ell,h}z_\ell-BF_{\ell,h}
\end{equation}
for any integral carry with $z_{n+1}=b_nz_n-Bm_n$; the integral form is the
\lword{Erdos269/RestrictedFloorSum.lean}{480}{integralCarry_window}{checked
window identity}, over the
\lword{Erdos269/RestrictedFloorSum.lean}{455}{affineRecurrence_window}{affine
window identity} and the
\lword{Erdos269/RestrictedFloorSum.lean}{469}{windowForcing_const_mul}{scaled
forcing identity}, and the real form is
\lword{Erdos269/CofinalWindowEscapeEquivalence.lean}{68}{trueNormalizedState_window}{checked}
for the genuine tail.  For $C>0$ let $\lpr_C(N)$ be the unique integer in
$\{1,\ldots,C\}$ congruent to $N$ modulo $C$, so $\lpr_C(0)=C$; this is the
\lword{Erdos269/ResidueEscape.lean}{26}{leastPositiveResidue}{canonical
representative}, with its
\lword{Erdos269/ResidueEscape.lean}{31}{leastPositiveResidue_pos_le}{range} and
its
\lword{Erdos269/ResidueEscape.lean}{52}{leastPositiveResidue_modEq}{congruence}
checked.  The vanishing residue is represented by $C$, and that choice is what
makes the next statement a modular one.

\begin{proposition}[the finite residue contradiction]\label{long269:res:consumer}
Let $C>0$ and let $c$ be an integer with $0<c$ and $|c|\le K$.  If
$c\equiv N\pmod C$ and $K<\lpr_C(N)$, then the hypotheses are contradictory.
\end{proposition}

\begin{proof}
The canonical representative lies in $\{1,\ldots,C\}$, and the inequalities put
$c$ strictly between $0$ and $C$.  If $N\equiv0\pmod C$ its representative is
$C$ while $c\bmod C=c\ne0$.  Otherwise $c$ and $\lpr_C(N)$ are each their own
residue and congruence makes them equal, contradicting $|c|\le K<\lpr_C(N)$.
\end{proof}

\noindent The abstract one-window predicate is the
\lword{Erdos269/ResidueEscape.lean}{71}{ResidueEscapesWindow}{residue-escape
condition}, with the contrapositive
\lword{Erdos269/ResidueEscape.lean}{96}{residue_le_bound_of_bounded_positive_state}{bounding
the residue from a bounded positive state}.  The natural-state version is
\lword{Erdos269/ResidueEscape.lean}{76}{no_bounded_positive_state_of_residue_escape}{checked},
the integer carry version is
\lword{Erdos269/ResidueEscape.lean}{110}{no_bounded_positive_int_state_of_leastPositiveResidue}{checked},
and the exact classifier $\lpr_C(N)=|c|$ under $|c|\le C$ is
\lword{Erdos269/ResidueEscape.lean}{138}{leastPositiveResidue_eq_natAbs_of_pos_le_modEq}{checked}
as well.  Applying it inside a window uses the forcing of the endpoint carry to
the canonical residue,
\lword{Erdos269/RestrictedFloorSum.lean}{497}{leastPositiveResidue_windowForcing_eq_carry}{checked}.

For a bound $G:\N_{>0}\times\N\to\N$ let $\Esc(G)$ be the statement
\begin{equation}\label{long269:eq:actual-escape}
 \begin{gathered}
 \text{for every }B\ge1\text{ with }\gcd(B,30)=1\text{ and every }a_0\ge1,\\
 \text{there are }\ell\ge a_0\text{ and }h\ge1\text{ with }
 \lpr_{W_{\ell,h}}(-BF_{\ell,h})>G(B,\ell+h) .
 \end{gathered}
\end{equation}
The window may depend on both $B$ and $a_0$.  The general predicate is the
\lword{Erdos269/RestrictedFloorSum.lean}{629}{CofinalLocalWindowEscape}{cofinal
local-window escape condition}, and the instance at the fixed bound
$K_0(B,a)=90B(a+1)^{2}$ is the
\lword{Erdos269/RationalityCarryBridge.lean}{397}{ActualCofinalLocalWindowEscape}{actual
producer}.  Every $b_a$ is positive, so $W_{\ell,h}>0$ is automatic.

\begin{theorem}[the equivalence band]\label{long269:res:actual-escape-endpoint}
\label{long269:res:lead-escape-equivalence}\label{long269:res:windowconsumer}
Let $G:\N_{>0}\times\N\to\N$ satisfy $K(B,a)\le G(B,a)$ for all $B$ and $a$,
and $G(B,a)/2^{a}\to0$ as $a\to\infty$ for each fixed $B$.  Then
\[
 \Esc(G)\quad\Longleftrightarrow\quad S\notin\Q .
\]
Both $K$ of~\eqref{long269:eq:actual-bound} and $K_0(B,a)=90B(a+1)^{2}$ lie in this
band, so $\Esc(K)$, $\Esc(K_0)$ and irrationality of $S$ are mutually
equivalent.  The lower half of the band cannot be dropped: $\Esc(0)$ holds
vacuously, since every least positive residue is at least $1$.
\end{theorem}

\begin{proof}
Suppose $\Esc(G)$ and suppose $S=N/D$ were rational.
Theorem~\ref{long269:res:actual-cancellation} supplies $B\ge1$ coprime to $30$, an
onset $a_D$, and positive integers $z_a=BX_a\le K(B,a)\le G(B,a)$ for
$a\ge a_D$ satisfying the cleared recurrence.  Apply~\eqref{long269:eq:actual-escape}
with $a_0=a_D$ to obtain a window $(\ell,h)$ with $\ell\ge a_D$.
By~\eqref{long269:eq:window-identity}, $z_{\ell+h}\equiv-BF_{\ell,h}$ modulo
$W_{\ell,h}$, while $0<z_{\ell+h}\le G(B,\ell+h)<\lpr_{W_{\ell,h}}(-BF_{\ell,h})$.
Proposition~\ref{long269:res:consumer} is the contradiction, so $S$ is irrational.

Conversely suppose $S\notin\Q$, and fix $B\ge1$ coprime to $30$ and $a_0\ge1$.
Set $\ell=a_0$.  In $X_\ell=h_\ell\bigl(S-\sum_{a<\ell}s_a\bigr)$ the finite
prefix is rational and $h_\ell$ is a positive rational, so $BX_\ell$ is
irrational and its distance $\delta$ from $\Z$ is positive.  Suppose no length
$h$ escaped, so that $r_h=\lpr_{W_{\ell,h}}(-BF_{\ell,h})\le G(B,\ell+h)$ for
every $h\ge1$.  Then $k_h=(BF_{\ell,h}+r_h)/W_{\ell,h}$ is an integer, and
multiplying~\eqref{long269:eq:window-identity} by $B$ gives
\[
 BX_\ell-k_h=\frac{BX_{\ell+h}-r_h}{W_{\ell,h}} .
\]
Both $BX_{\ell+h}$ and $r_h$ are positive, the first at most
$BQ(n_{\ell+h})$ by Theorem~\ref{long269:res:actual-tail-bound} and the second at most
$G(B,\ell+h)$, so the numerator has absolute value at most
$\max\bigl(BQ(n_{\ell+h}),G(B,\ell+h)\bigr)$.  Each $b_a\ge2$ gives
$W_{\ell,h}\ge2^{h}$, so
\[
 0<\delta\le|BX_\ell-k_h|
 \le\frac{\max\bigl(BQ(n_{\ell+h}),\,G(B,\ell+h)\bigr)}{2^{h}}
 =2^{\ell}\cdot
 \frac{\max\bigl(BQ(n_{\ell+h}),\,G(B,\ell+h)\bigr)}{2^{\ell+h}}
 \longrightarrow0
\]
as $h\to\infty$, because $Q(n_{\ell+h})$ is quadratic in $\ell+h$ and
$G(B,a)=o(2^{a})$.  This contradiction proves $\Esc(G)$.

For the two named bounds, $K\le K_0$ by
Theorem~\ref{long269:res:actual-cancellation}, both are quadratic in $a$ for fixed $B$,
and $K\le K$ trivially.  Finally $\lpr_C(N)\ge1>0$ always, so $\Esc(0)$ holds
whatever $S$ is, while $0\le K$ fails.
\end{proof}

\noindent The converse direction is checked in Lean at a strictly greater
generality than the fixed bound: irrationality implies escape against any short
bound that the window growth eventually beats,
\lword{Erdos269/CofinalWindowEscapeEquivalence.lean}{327}{cofinalLocalWindowEscape_of_irrational_of_beaten}{checked},
and in particular against every bound dominated by $c(B)(n+1)^{2}$,
\lword{Erdos269/CofinalWindowEscapeEquivalence.lean}{355}{cofinalLocalWindowEscape_of_irrational_of_quadratic}{checked}.
The forward direction at the fixed bound $K_0$ is
\lword{Erdos269/RationalityCarryBridge.lean}{479}{irrational_value_of_cofinalLocalWindowEscape}{checked},
over the abstract consumer
\lword{Erdos269/RestrictedFloorSum.lean}{645}{no_positive_reducedCarry_of_cofinalLocalWindowEscape}{checked}
and its packaged absorbed form
\lword{Erdos269/RestrictedFloorSum.lean}{689}{no_positive_absorbedCarry_of_cofinalLocalWindowEscape}{checked}.
The two directions combine into the equivalence
\lword{Erdos269/CofinalWindowEscapeEquivalence.lean}{392}{actualCofinalLocalWindowEscape_iff}{checked}
for the shifted tail and
\lword{Erdos269/CofinalWindowEscapeEquivalence.lean}{398}{actualCofinalLocalWindowEscape_iff_irrational_value}{checked}
for the series value itself.  The Lean equivalence is stated at $K_0$; the band
of Theorem~\ref{long269:res:actual-escape-endpoint} is the paper statement.

Theorem~\ref{long269:res:actual-escape-endpoint} settles a question that a reader is
entitled to ask about the criterion.  Sharpening the bound below $K$ does not
weaken the producer, since escape against a bound that fails to dominate the
actual carries proves nothing; the zero bound is the extreme case.  Enlarging
the bound up to any subexponential function does not weaken it either.  The
producer is therefore a restatement of Erd\H{o}s~\#269 for $P=\{2,3,5\}$
throughout the band, and work on it is work on the target.  Equivalence
preserves truth and settles nothing about difficulty, so the reformulation may
still be the easier representation to attack; what it does not admit is a
cheaper bound.

Two exact countermodels rule out further shortcuts.  For $(W,F,B)=(6,4,1)$ we
have $\lpr_6(-4)=2$, so the canonical residue need not be coprime to the
accumulated base.  For $(W,F,B)=(60,47,37)$ we have
$\lpr_{60}(-37\cdot47)=1$, so a fixed window has no denominator-independent
lower bound on that residue.  The window may therefore depend genuinely on $B$.

\subsection{How fast the window base grows}

The crude estimate $W_{\ell,h}\ge2^{h}$ is all the equivalence needs.  The
actual height formula gives more, and the extra information sets the scale of
any finite search.

\begin{proposition}[window-growth law]\label{long269:res:window-growth}
Put $\theta_3=\log_32$ and $\theta_5=\log_52$.  For all $\ell\ge0$ and
$h\ge1$,
\[
 W_{\ell,h}=2^{h}\,
 3^{\lfloor(\ell+h)\theta_3\rfloor-\lfloor\ell\theta_3\rfloor}\,
 5^{\lfloor(\ell+h)\theta_5\rfloor-\lfloor\ell\theta_5\rfloor},
 \qquad
 \frac{8^{h}}{15}<W_{\ell,h}<15\cdot8^{h} .
\]
\end{proposition}

\begin{proof}
Telescoping \eqref{long269:eq:dyadic-alphabet} gives
$W_{\ell,h}=\hgt_{\ell+h}/\hgt_\ell$, and the displayed formula is that
quotient written out.  Each floor difference differs from $h\theta_p$ by less
than one, and $3^{\theta_3}=5^{\theta_5}=2$, so the $3$-factor lies strictly
between $2^{h}/3$ and $3\cdot2^{h}$ and the $5$-factor strictly between
$2^{h}/5$ and $5\cdot2^{h}$.  Multiplying the three ranges gives the bounds.
\end{proof}

\begin{corollary}[no bounded-length escape at all starts]\label{long269:res:no-bounded-length}
Fix $B\ge1$ coprime to $30$ and $H\ge1$.  Only finitely many starts $\ell$
admit an escaping window of length at most $H$ against the bound $K$.
\end{corollary}

\begin{proof}
A least positive residue never exceeds its modulus, so escape at $(\ell,h)$
requires $K(B,\ell+h)<W_{\ell,h}<15\cdot8^{h}\le15\cdot8^{H}$.  On the other
hand $j_a\ge a-1$, since the powers $2,\ldots,2^{a-1}$ already lie below
$2^{a}$, so $K(B,\ell+h)\ge\lfloor B((\ell-1)^{2}+10(\ell-1)+27)/9\rfloor$,
which tends to infinity with $\ell$.
\end{proof}

\noindent Corollary~\ref{long269:res:no-bounded-length} is the exact reason a finite
scan cannot approach the cofinal quantifier by widening its denominator range
alone.  Rearranging its inequality, an escaping window at start $\ell$ and
endpoint $a=\ell+h$ must satisfy
\[
 3h>\log_2K(B,a)-\log_215,
\]
so the necessary search depth grows like
$\bigl(\log_2B+2\log_2\ell\bigr)/3$.  This is a lower bound on the length that can
possibly work, and it is not an upper bound on the first length that does.
Producing the latter is exactly the open problem, and by
Theorem~\ref{long269:res:actual-escape-endpoint} it needs effective control of
$\operatorname{dist}(BX_\ell,\Z)$: counting window growth is easy, and keeping
a scaled tail away from the integers is the arithmetic content.

\statusboundary

\section{The finite evidence}\label{long269:sec:evidence}

Three finite computations bear on the problem.  Each is exact integer or
certified-enclosure arithmetic, none is a Lean theorem, and none settles an
instance.

\subsection{The dyadic-window scan}

A \emph{certificate} is a finite tuple of integers recording one instance of
the inequality in~\eqref{long269:eq:actual-escape}, so that a reader can recheck it in
a line.  The integer-only
\href{\sourceurl/scripts/check_erdos269_dyadic_windows.py}{dyadic-window
checker} constructs the ordered pure-power jumps, the block bases and the block
digits from exact multiplicity counts, and reproduces the following, with
columns the denominator $B$, the start $\ell$, the length $h$, the endpoint
jump index $j_{\ell+h}$, the window base, the forcing, the residue
$R=\lpr_{W}(-BF)$ and the bound $K$ of~\eqref{long269:eq:actual-bound}.
\[
\begin{array}{c|c|c|c|r|r|r|r}
B&\ell&h&j_{\ell+h}&W&F&R&K\\ \hline
1&1&2&4&60&47&13&9\\
7&1&3&5&360&289&137&95\\
16&1&4&7&10800&8735&640&352
\end{array}
\]
The first row reads as follows.  The window starts at $\ell=1$ and has length
$2$, so $W=b_1b_2=6\cdot10=60$; the accumulated forcing is $F=47$; and
$\lpr_{60}(-47)=13$, since $-47+60=13$, which exceeds
$K(1,3)=\lfloor(16+40+27)/9\rfloor=9$.  The third row lies outside the domain
of~\eqref{long269:eq:actual-escape}, since $\gcd(16,30)=2$, and is displayed to
illustrate the window arithmetic at greater depth.

A fresh scan over every $B\le5000$ coprime to $30$ and every $100\le\ell\le3000$
tested $3{,}869{,}934$ pairs and found an escaping window in every case, of
length at most $18$ with search depth permitted to $24$.  The distribution of
first successful lengths concentrates at $10$ to $12$, which is where
Corollary~\ref{long269:res:no-bounded-length} predicts the shortest possible window to
lie over this range.  The computation uses integers only and is reproducible
from the pinned checker with
\texttt{-{}-max-denominator 5000 -{}-start-min 100 -{}-start-max 3000
-{}-max-length 24 -{}-assert-packet}.  Neither the scan nor the three displayed
certificates proves escape for unbounded $B$ or for cofinally many starts, and
by Corollary~\ref{long269:res:no-bounded-length} no scan at bounded length can.

\subsection{Two finite denominator exclusions}

\begin{theorem}[window-$128$ block exclusion, computational certificate]
\label{long269:res:lead-block-exclusion}
At window length $L=128$, launch $a_1=10005$ and $64$ starts, no rational value
of the $\{2,3,5\}$ running-LCM series has reduced denominator $MB$ with $M$ a
$30$-smooth divisor of $2^{10005}3^{6312}5^{4308}$, $\gcd(B,30)=1$ and
$1<B\le B_{\max}$, where $B_{\max}$ is the $106$-digit integer
\[
\begin{aligned}
 B_{\max}={}&1134599670999687767349520845707093359257353022286558739363600235\\
 &016103207564063373270305324172145281971729 ,
\end{aligned}
\]
so that $\log_2B_{\max}=348.9846\ldots$
\end{theorem}

The exponent triple $(10005,6312,4308)$ is the height exponent triple of
$\hgt(2^{10005})$, so the certificate normalises by the full height at its
launch.  Lemma~\ref{long269:res:all-scale-lattice} normalises by the half height
$h_{10005}$; the certificate's smooth family is accordingly the larger one.  The recorded quantities are the window product with
$\log_2P=386.40993\ldots$, the exclusion index $J=1$ certified from a reduced
basis lying inside the per-start budget, an enclosure of width
$9.674\times10^{-227}$, and a maximum ratio $\max X/W=0.185997$.  This is a
computational certificate and nothing else: no Lean declaration carries any
part of it, so it rests on the correctness of the exact integer engine that
produced it, and the certificate record is held in the author's
formal-mathematics archive, outside the public source of this release.

\begin{theorem}[continued-fraction exclusion, computational certificate]
\label{long269:res:cf-exclusion}
The normalised tail $X_1$ has $13{,}109$ certified partial quotients, so it is
not rational with denominator at most $2^{22482}$, about $10^{6768}$.
\end{theorem}

The certification is by common prefix of the continued fractions of the two
endpoints of an interval provably containing $X_1$, so no approximation
heuristic enters, and the truncation was independently checked to agree with
the direct smooth-number sum as an exact rational.  The recorded statistics are
a largest denominator of $22{,}483$ bits, a largest partial quotient of
$129{,}114$, a mean partial quotient of $23.4133$, observed Gauss--Kuzmin
frequencies $0.4208$, $0.1665$, $0.0917$, $0.0575$, $0.0391$ against the
predicted $0.4150$, $0.1699$, $0.0931$, $0.0589$, $0.0406$, and a L\'evy
constant of $1.18869$ against $\pi^{2}/(12\log2)=1.18657$.  No Liouville
behaviour and no algebraic or self-similar continued-fraction structure appears
below that height.

Each exclusion is finite and neither contains the other: one is indexed by a
lattice at a fixed launch and fixed smooth part, the other by continued-fraction
depth at $a=1$.  Every larger denominator survives both, so neither settles an
instance of the problem.

\section{The remaining arithmetic questions}\label{long269:sec:open}

\begin{problem}[actual cofinal residue escape]\label{long269:prob:producer}
Prove $\Esc(K)$ for the actual digits~\eqref{long269:eq:actual-digit} and the bound
\eqref{long269:eq:actual-bound}.
\end{problem}

Every term in this question is a finite integer quantity, and a proof must
produce a later window for every reduced denominator and every prescribed
onset.  The tail estimate and the smooth-factor cancellation are proved above
and are not additional hypotheses.  By
Theorem~\ref{long269:res:actual-escape-endpoint} the problem is a restatement of
Erd\H{o}s~\#269 for $P=\{2,3,5\}$.

\begin{problem}[nonintegrality of every reduced tail]\label{long269:prob:tails269}
For every $B\ge1$ with $\gcd(B,30)=1$ and every $a\ge1$, prove
\begin{equation}\label{long269:eq:tail-nonintegrality}
 BX_a\notin\Z .
\end{equation}
\end{problem}

\begin{proposition}[the reduced-tail question is also the target]
\label{long269:res:tails-equivalence}
Statement~\eqref{long269:eq:tail-nonintegrality}, quantified over every $B\ge1$
coprime to $30$ and every $a\ge1$, is equivalent to irrationality of $S$.
\end{proposition}

\begin{proof}
If $BX_a\in\Z$ for some such $B$ and $a$, then $X_a\in\Q$, and since
$S=\sum_{j<a}s_j+X_a/h_a$ with $h_a$ a positive rational and the prefix a
finite sum of rationals, $S\in\Q$.  Conversely if $S\in\Q$, then
Theorem~\ref{long269:res:actual-cancellation} produces $B$ coprime to $30$ with
$BX_a\in\Z$ for every $a\ge a_D$.
\end{proof}

\noindent The abstract form of the first implication is
\lword{Erdos269/BoundedRadixTailEscape.lean}{183}{rational_of_scaledTail_integer}{checked}.
Problem~\ref{long269:prob:tails269} is therefore a second restatement of the target,
and its pointwise shape is the more convenient one: if $BX_a$ is integral at one index, the integer-coefficient recurrence
makes it integral at every later index, so a single index decides it.
Proposition~\ref{long269:res:actual-dichotomy} does not exclude that branch, since it
constrains $X_a$ alone and a rational value may keep a surviving denominator
coprime to $30$.

\subsection{The analytic route and what it would need}

For the de-duplicated series put
\[
 \mathcal D_{2,3,5}=1+
 \sum_{t\in\{2^{n},3^{n},5^{n}:n\ge1\}}
 \frac{1}{2^{\lfloor\log_2t\rfloor}3^{\lfloor\log_3t\rfloor}
           5^{\lfloor\log_5t\rfloor}} ,
\]
whose dyadic coding is the joint rotation word
$\delta_{3,a}=\lfloor(a+1)\theta_3\rfloor-\lfloor a\theta_3\rfloor$ and
$\delta_{5,a}=\lfloor(a+1)\theta_5\rfloor-\lfloor a\theta_5\rfloor$, with
$b_a=2\cdot3^{\delta_{3,a}}5^{\delta_{5,a}}$.

\begin{problem}[function-faithful two-dimensional representation]
\label{long269:prob:representation}
Express $\mathcal D_{2,3,5}$ as a nonconstant algebraic combination of values
of a specified two-dimensional Hecke--Mahler, cone-generating or multivariate
Mahler function and verify every hypothesis of a published value theorem; or
give a conditional theorem under an explicit logarithmic nondegeneracy
hypothesis; or prove that the literal series has no representation in the
specified finite-dimensional class.
\end{problem}

Individual irrationality of $\theta_3$ and $\theta_5$ is what
Theorem~\ref{long269:res:infinite-rank} uses, and it is weaker than the joint
equidistribution or the rational independence of $1,\theta_3,\theta_5$ that a
two-dimensional value theorem may need; the stronger hypothesis is not assumed
anywhere above.  The finite-observer formalisation isolates the precise
faithfulness requirement: equality in a finite observer must imply equality
after symbolic realisation, and a genuine finite-dimensional factorisation
forces the realised symbolic span to be finite-dimensional
(\lword{Erdos269/WeightedPhaseCarry.lean}{109}{carry_eq_residueDigit_add_coboundary}{residue-coboundary
form},
\lword{Erdos269/WeightedPhaseCarry.lean}{293}{CartierRealisation}{symbolic
realisation},
\lword{Erdos269/WeightedPhaseCarry.lean}{334}{finite_realisedSpan_of_factorisation}{checked}),
with the carry residue and residue digit confined to their declared intervals
(\lword{Erdos269/WeightedPhaseCarry.lean}{150}{carryResidue_mem_interval}{carry
interval},
\lword{Erdos269/WeightedPhaseCarry.lean}{157}{residueDigit_mem_interval}{digit
interval}).  No theorem here proves that the literal realised span is infinite.

\subsection{Structural constraints already in place}

The radix alphabet extends without difficulty.  For ordered primes
$p_1<\cdots<p_s$ an interval $(p_1^{a},p_1^{a+1})$ contains at most one power
from each other channel, because consecutive $p_i$-powers have ratio
$p_i>p_1$, so its block radix belongs to the $2^{s-1}$-letter alphabet
$\{p_1\prod_{i=2}^{s}p_i^{\varepsilon_i}:\varepsilon_i\in\{0,1\}\}$.  The
quantitative questions are the interesting ones: effective recurrence or
discrepancy for the actual four-letter $\{2,6,10,30\}$ word, an asymptotic with
an error term for the restricted two-dimensional shell counts that generate
$m_a$, or the exact separated rank of the literal kernel under a specified
family of shifts.

Three-channel rigidity and carry-lift extinction already exclude one proposed
argument in four exact steps.  Under channel surjectivity, ordinary block
nullity is equivalent to the perturbation being a coboundary of a channel
potential
(\lword{Erdos269/ThreeChannelBlockRigidity.lean}{59}{channelBlockNull_iff_channelPotential}{potential
classifier}).  Zero perturbations on genuine $2\to3$ and $2\to5$ transitions
then identify all three potential values and force the perturbation to vanish
at every index.  For an integral lift, that vanishing makes a nonzero initial
error grow by the exact product of the successive bases, and bases at least two
make its absolute value at least $2^{N}$, contradicting even a single
index-$N$ bound strictly below $2^{N}$
(\lword{Erdos269/CarryLiftExtinction.lean}{178}{no_carryLift_of_errorBound_below_twoPow}{one-index
extinction}); consequently no uniform bound on the lift error can hold either
(\lword{Erdos269/CarryLiftExtinction.lean}{238}{no_bounded_carryLift_of_blockNull_twoAnchors}{uniform
extinction}).  A separate four-state calculation reaches the obstruction
earlier: four real states in $(0,1)$ with unit-accuracy integral lifts and the
two anchor equalities force the first complete $2$-block sum to be $1$, so that
block cannot be null
(\lword{Erdos269/CarryLiftExtinction.lean}{289}{binaryLift_twoAnchors_force_firstBlock_sum_one}{first-block
sum},
\lword{Erdos269/CarryLiftExtinction.lean}{308}{no_unitAccurateLift_with_twoAnchors_and_firstTwoBlockNull}{four-state
obstruction}).

These lift conclusions are conditional.  The paper constructs the actual orbit
and its bounded integral carry under rationality, and it does not construct a
faithful lift with the two anchors or the block nullity those auxiliary
statements require.  The actual carry supplies a weighted block defect instead,
so any successful argument along that line must use the weighted identity or
construct a different faithful lift.

A single conditional single-channel criterion is also on record.  For the pure
$2$-channel sum the Cantor states lie in an explicit open interval
(\lword{Erdos269/PurePowerIrrationality.lean}{61}{state_mem_Ioo}{confinement}),
consecutive states are separated
(\lword{Erdos269/PurePowerIrrationality.lean}{74}{stateGap_lower_bound}{gap
bound}), a small nonzero linear form forces irrationality
(\lword{Erdos269/PurePowerIrrationality.lean}{94}{irrational_of_small_nonzero_forms}{criterion}),
and clearing together with small gaps assembles to irrationality
(\lword{Erdos269/PurePowerIrrationality.lean}{156}{irrational_of_clearing_and_small_gaps}{assembly}).
The clearing and small-gap hypotheses are not discharged for any concrete
series, so no irrationality statement follows from them here.

\subsection{Where the problem stands}

Erd\H{o}s~\#269 is settled for $|P|=2$, at the level of transcendence, and open
for every $|P|\ge3$.  For $P=\{2,3,5\}$ the reduction is complete on the
arithmetic side: the literal digit, the tail, the integral recurrence, the
smooth-factor cancellation, the explicit onset and the sharp carry bound all
refer to the same series, and Theorem~\ref{long269:res:actual-escape-endpoint} shows
that the one remaining condition is the problem itself throughout an entire
band of bounds.  Three obstructions mark the limits.  Arbitrary-order minors
exclude every finite exact separation of the kernel, and
Theorem~\ref{long269:res:uniform-rank} excludes even approximate separation of its
normalised carry below half a jump.  The bounded-radix alternative leaves
integral tails untouched.  The residue countermodels show that coprimality
alone gives no denominator-independent residue bound, and
Corollary~\ref{long269:res:no-bounded-length} shows that no bounded-length search can
reach the cofinal quantifier.  What remains is to force the literal word's
residue outside the finite carry interval, with the window allowed to depend on
the denominator and to begin beyond the onset.

\pdfbookmark[1]{Statements and declarations}{statements-declarations}
\subsection*{Statements and declarations}

\paragraph{Evidence.}  Every linked phrase above opens its Lean declaration at
the pinned source revision~\commitshort, where the finite running-LCM geometry,
the two-prime outer product, the arbitrary-order kernel minors and the
exclusion of finite separation, the literal shell summability and recurrence,
the all-scale rationality lattice, the rationality-to-carry bridge, the residue
arithmetic and both directions of the window-escape equivalence are checked by
the Lean~4 kernel against Mathlib with no \texttt{sorry}.  The two-prime
transcendence theorem is an ordinary proof over the cited Bugeaud--Laurent
theorem and carries no Lean declaration; the quadratic tail bound, the explicit
onset $a_D$, the sharp bound $K$, the equivalence band, the window-growth law
and Theorem~\ref{long269:res:uniform-rank} are ordinary proofs in this paper; the three
finite computations of Section~\ref{long269:sec:evidence} are computational
certificates.

\paragraph{Artefact and data availability.} The
\href{\sourceurl}{pinned formal-source revision} contains the Lean sources, the
fixed toolchain, the library manifest, and the exact dyadic-window checker used
in the finite scan.  Proof authority rests in those pinned sources.  The
present text is exposition and navigation.

\paragraph{Funding and competing interests.} This work received no external
funding.  The author declares no competing interests.

\paragraph{Acknowledgements.} The two-prime factorisation, its Hecke--Mahler
reduction and the running-LCM identity for every $|P|$ are credited to Steve
Fan.  The transcendence input is credited to Yann Bugeaud and Michel Laurent
and to the earlier work of Loxton and van der Poorten cited by them.  The
problem numbering and status follow the Erd\H{o}s Problems catalogue maintained
by Thomas Bloom~\cite{erdosproblems}.
% ---- part extended_record ----
\section{Extended record: migrated arguments and finite evidence}
\label{long269:long:extended}

This part of the record holds material that the short note compressed or moved.
Nothing here is an additional result; every statement is either an auxiliary
form of a theorem of the note, a historical or bibliographic detail, or the
exact source inventory.  The canonical references are the note's own labels.

\subsection{The historical and catalogue record}\label{long269:long:history}

In the primary 1988 source Erd\H{o}s states the infinite-$P$ assertion as a
simple exercise and presents persistence for a finite number of primes greater
than one as a probable extension, and not as a theorem
\cite[p.~106]{erdos1988}.  The current open status is supplied by the
catalogue \cite{erdosproblems}, and is not inferred from that conjectural
wording.  The 1974 letter writes ``given primes $p_1,\ldots,p_r$'' without
explicitly restricting $r$, so the singleton case, whose value is $p/(p-1)$,
must be excluded by hand; the modern restriction $|P|\ge2$ does that.

The letter's de-duplicated assertion is made for a general finite list of
primes and supplies no proof, so it is an asserted historical result and not an
open problem.  The note therefore restricts its open statement to the repeated
series $\mathcal R_P$.  A recovered proof of the letter's assertion, or a
transcendence statement for $\mathcal D_P$ with $|P|\ge3$, would be a different
question and needs its own formulation.

The two-prime argument of the note is independent of the unprinted argument in
the letter, and it is not the first public proof.  Steve Fan posted the same
factorisation, the same Hecke--Mahler reduction and the same conclusion in the
discussion thread of the problem's page on 26 June 2026
\cite{fan2026comment}, with follow-up remarks extending the argument to
arbitrary coprime pairs; a further comment there observes that the argument
does not seem to generalise immediately to $|P|\ge3$.  The manuscript of the
note was first released publicly on 22 July 2026, at commit \texttt{a9d3ab8}.

The statement of the problem has been formalised as a conjecture with an
unfilled proof in the \emph{Formal Conjectures} collection
\cite{formalconjectures269}.  Its Nat-indexed series includes the empty-prefix
least-common-multiple term, so its value is exactly \(1\) plus the conventional
series; transporting a theorem across that boundary needs an
explicit series-identification lemma, which is not supplied here.

\subsection{The abstract carry lemmas in their general form}
\label{long269:long:general-carries}

The note applies the carry machinery to the literal $\{2,3,5\}$ orbit.  The
underlying statements are about integer sequences alone, they are reusable, and
they are recorded here in the generality in which they are checked.

Let $D=D_{\mathrm{sm}}B$ with $D_{\mathrm{sm}}=2^{u}3^{v}5^{w}$ and
$\gcd(B,30)=1$, and let $(c_n)$ be an integer sequence satisfying
$c_{n+1}=b_nc_n-Dm_n$ for an integer radix word $(b_n)$ and forcing word
$(m_n)$.

\begin{proposition}[conditional denominator reduction]
\label{long269:long:denominator-reduction}
If $c_n=D_{\mathrm{sm}}d_n$ for every $n$, with $D_{\mathrm{sm}}>0$, then the
recurrence, positivity, upper bound and window identity for $(c_n)$ reduce to
the same four statements for $(d_n)$ with multiplier $B$ in place of $D$.
\end{proposition}

Height absorption alone does not imply divisibility of a carry state.  The
identity $DX_a=h_aN-Dv_a$ of Lemma~\ref{long269:res:all-scale-lattice} is what supplies
it for the actual orbit, and Theorem~\ref{long269:res:actual-cancellation} is the
instance in which the note uses it.  The formal consumer takes the
common-factor form as a hypothesis.

\begin{proposition}[conditional extinction of bounded carries]
\label{long269:long:windowconsumer}
Let $(b_n)$ and $(m_n)$ be a radix word and a forcing word, let
$G:\N_{>0}\times\N\to\N$, and assume cofinal local-window escape against $G$.
Fix $B>0$ coprime to $30$.  There is no integral sequence $(d_n)$ satisfying
simultaneously $d_{n+1}=b_nd_n-Bm_n$, $d_n>0$ and $|d_n|\le G(B,n)$ for every
$n\ge0$.
\end{proposition}

\begin{proof}
Choose one escaping window $(\ell,h)$.  The window identity gives
$d_{\ell+h}\equiv-BF_{\ell,h}$ modulo $|W_{\ell,h}|$.  The endpoint state is
positive and at most $G(B,\ell+h)$, whereas the canonical positive residue of
the right-hand side exceeds that bound, so Proposition~\ref{long269:res:consumer}
applies.
\end{proof}

Coprimality with $30$ is used by the escape hypothesis to select a window; once
a window is fixed, the finite contradiction does not use it.  The formalisation
carries the edge cases: a zero window base is excluded, a zero residue is
represented by the full modulus, and positivity prevents the endpoint carry
from vanishing.  The packaged absorbed form takes a nonzero smooth factor, the
exact factorisation $c_n=D_{\mathrm{sm}}d_n$, the absorbed recurrence for
$(c_n)$ and the positive short bound for $(d_n)$, and derives the same
contradiction in one statement.

\subsection{Worked finite examples}\label{long269:long:finite-geometry}

The ten smallest values of the running least common multiple at $\{2,3,5\}$ are
tabulated in Section~\ref{long269:sec:lcm}.  They illustrate both parts of
Proposition~\ref{long269:res:cell}: the value is constant on $\{5,6,7\}$ and on
$\{9,10\}$, and each change multiplies by a single prime, by $2$ at
$x=2,4,8$, by $3$ at $x=3,9$ and by $5$ at $x=5$.  Also $\Lprefix(10)=8\cdot9\cdot5=360$
is the least common multiple of the smooth numbers $1,2,3,4,5,6,8,9,10$.

For Proposition~\ref{long269:res:fibre} take $(p,q,r)=(2,3,5)$ and the box
$\mathcal B(1,1,1)$, whose eight points carry the smooth values
$1,2,3,5,6,10,15,30$ and the heights
\[
\begin{array}{c|cccccccc}
p^{i}q^{j}r^{k}&1&2&3&5&6&10&15&30\\ \hline
\hgt&1&2&6&60&60&360&360&10800
\end{array}
\]
Six heights occur, two of them twice: the points $5$ and $6$ share the height
$60$, and $10$ and $15$ share the height $360$.  The identity reads
\[
 1+\tfrac12+\tfrac16+\tfrac1{60}+\tfrac1{60}+\tfrac1{360}+\tfrac1{360}
 +\tfrac1{10800}
 =1+\tfrac12+\tfrac16+\tfrac2{60}+\tfrac2{360}+\tfrac1{10800}
 =\tfrac{18421}{10800},
\]
and the two coefficients $2$ carry the whole content of the regrouping on this
box.  Where the heights are pairwise distinct the identity is a relabelling.

Multiplying block radices along a run of blocks gives the product of the
jump-word letters over that run: for instance $b_1b_2=60$ is the product of the
four multipliers at $3,4,5,8$.  Block $4$ is the only one among the first six
carrying an internal power in both channels, and block $5$ carries neither, so
its radix falls back to the terminal factor alone.

\subsection{The finite-scan histogram}\label{long269:long:experiments}

The scan of Section~\ref{long269:sec:evidence} covers every $B\le5000$ coprime to $30$
and every start $100\le\ell\le3000$, with search depth permitted to $24$.  It
tested $3{,}869{,}934$ pairs, found an escaping window in every case, and the
distribution of first successful lengths was
\[
\resizebox{\linewidth}{!}{$\begin{array}{c|rrrrrrrrrrrrrrr}
h&4&5&6&7&8&9&10&11&12&13&14&15&16&17&18\\ \hline
\#&1&104&812&5437&51409&237423&735450&1431226&1132756&236752&34910&3076&521&49&8
\end{array}$}
\]
The first case at the maximal observed length $18$ is $B=917$ at start
$\ell=2980$, with endpoint jump index $6179$, window base
$18139852800000000$, forcing $13196471407660025821045$, residue
$76322101735$ and bound $3896420420$.  An earlier run of the same checker over
every $B\le1000$ coprime to $30$ and every start $100\le\ell\le500$ tested
$106{,}666$ pairs with maximal first successful length $14$, whose first case
is $B=359$ at start $291$, endpoint jump index $627$, base $5038848000000$,
forcing $25864575212865807$, residue $213175287$ and bound $15932659$.

By Corollary~\ref{long269:res:no-bounded-length} these histograms describe a bounded
region.  The observed concentration at lengths $10$ to $12$ sits just above the
necessary depth $\bigl(\log_2K(B,\ell+h)-\log_215\bigr)/3$ supplied by
Proposition~\ref{long269:res:window-growth}, which is what a reader should expect if
the residues behave like generic ones inside this range.

\subsection{Source inventory}\label{long269:long:sources}

Each entry names an informal statement and the Lean declaration that carries it
at the pinned source revision~\commitshort.

\begin{longtable}{@{}p{.35\linewidth}p{.59\linewidth}@{}}
informal statement & linked Lean source\\
\endhead
smooth lattice value & \lref{Erdos269/ThreePrimeRunningLcm.lean}{32}{smooth3Val}\\
pure-power height & \lref{Erdos269/ThreePrimeRunningLcm.lean}{37}{threePrimeHeight}\\
lattice kernel & \lref{Erdos269/ThreePrimeRunningLcm.lean}{41}{threePrimeKernelQ}\\
smooth prefix index set & \lref{Erdos269/ThreePrimeRunningLcm.lean}{47}{smoothPrefixExponents}\\
running least common multiple & \lref{Erdos269/ThreePrimeRunningLcm.lean}{54}{smoothPrefixLcm}\\
prefix value divides the height & \lref{Erdos269/ThreePrimeRunningLcm.lean}{60}{smooth3Val_dvd_threePrimeHeight_of_mem}\\
divisibility into the height & \lref{Erdos269/ThreePrimeRunningLcm.lean}{78}{smoothPrefixLcm_dvd_threePrimeHeight}\\
first pure-power membership & \lref{Erdos269/ThreePrimeRunningLcm.lean}{85}{pureFirst_mem_smoothPrefixExponents}\\
second pure-power membership & \lref{Erdos269/ThreePrimeRunningLcm.lean}{98}{pureSecond_mem_smoothPrefixExponents}\\
third pure-power membership & \lref{Erdos269/ThreePrimeRunningLcm.lean}{111}{pureThird_mem_smoothPrefixExponents}\\
running-lcm identity & \lref{Erdos269/ThreePrimeRunningLcm.lean}{124}{smoothPrefixLcm_eq_threePrimeHeight}\\
cell relation & \lref{Erdos269/ThreePrimeRunningLcm.lean}{164}{SameThreePrimeLogCell}\\
cell constancy of the height & \lref{Erdos269/ThreePrimeRunningLcm.lean}{170}{threePrimeHeight_eq_of_sameLogCell}\\
cell constancy of the running value & \lref{Erdos269/ThreePrimeRunningLcm.lean}{180}{smoothPrefixLcm_eq_of_sameLogCell}\\
cell constancy of the kernel & \lref{Erdos269/ThreePrimeRunningLcm.lean}{192}{threePrimeKernelQ_eq_of_sameLogCell}\\
positive power sets & \lref{Erdos269/ThreePrimeRunningLcm.lean}{205}{positivePrimePowers}\\
channel cardinality & \lref{Erdos269/ThreePrimeRunningLcm.lean}{209}{positivePrimePowers_card}\\
exclusion of the origin & \lref{Erdos269/ThreePrimeRunningLcm.lean}{220}{one_not_mem_positivePrimePowers}\\
channel disjointness & \lref{Erdos269/ThreePrimeRunningLcm.lean}{230}{positivePrimePowers_disjoint}\\
positive jump channels & \lref{Erdos269/ThreePrimeRunningLcm.lean}{244}{threePrimePositiveJumpSet}\\
positive jump count & \lref{Erdos269/ThreePrimeRunningLcm.lean}{250}{threePrimePositiveJumpSet_card}\\
jump set with the origin & \lref{Erdos269/ThreePrimeRunningLcm.lean}{274}{threePrimeJumpSetWithOrigin}\\
jump count with the origin & \lref{Erdos269/ThreePrimeRunningLcm.lean}{279}{threePrimeJumpSetWithOrigin_card}\\
first height step & \lref{Erdos269/ThreePrimeRunningLcm.lean}{295}{threePrimeHeight_firstLogStep}\\
second height step & \lref{Erdos269/ThreePrimeRunningLcm.lean}{306}{threePrimeHeight_secondLogStep}\\
third height step & \lref{Erdos269/ThreePrimeRunningLcm.lean}{317}{threePrimeHeight_thirdLogStep}\\
first coordinate step & \lref{Erdos269/ThreePrimeRunningLcm.lean}{327}{smoothPrefixLcm_firstLogStep}\\
second coordinate step & \lref{Erdos269/ThreePrimeRunningLcm.lean}{340}{smoothPrefixLcm_secondLogStep}\\
third coordinate step & \lref{Erdos269/ThreePrimeRunningLcm.lean}{353}{smoothPrefixLcm_thirdLogStep}\\
exponent box & \lref{Erdos269/ThreePrimeRunningLcm.lean}{369}{smoothExponentBox}\\
point height & \lref{Erdos269/ThreePrimeRunningLcm.lean}{374}{smoothPointHeight}\\
height fibre & \lref{Erdos269/ThreePrimeRunningLcm.lean}{378}{smoothHeightFiber}\\
fibre sum & \lref{Erdos269/ThreePrimeRunningLcm.lean}{385}{smoothHeightFiber_kernel_sum}\\
height-fibre normal form & \lref{Erdos269/ThreePrimeRunningLcm.lean}{407}{finiteSmoothKernelSum_groupedByHeight}\\
cubic majorant & \lref{Erdos269/ThreePrimeRunningLcm.lean}{431}{threePrimeHeight_le_cube}\\
origin value & \lref{Erdos269/ThreePrimeRunningLcm.lean}{444}{kernel_235_origin}\\
value at two & \lref{Erdos269/ThreePrimeRunningLcm.lean}{451}{kernel_235_two}\\
value at three & \lref{Erdos269/ThreePrimeRunningLcm.lean}{459}{kernel_235_three}\\
value at six & \lref{Erdos269/ThreePrimeRunningLcm.lean}{468}{kernel_235_six}\\
non-separation witness & \lref{Erdos269/ThreePrimeRunningLcm.lean}{480}{kernel_235_not_rankOne}\\
variable-base tail step & \lref{Erdos269/ThreePrimeRunningLcm.lean}{488}{tailStateStep}\\
expanded tail step & \lref{Erdos269/ThreePrimeRunningLcm.lean}{491}{tailStateStep_eq}\\
smooth exponent shell & \lref{Erdos269/ThreePrimeRunningLcm.lean}{501}{smoothExponentShell}\\
short-interval uniqueness & \lref{Erdos269/ThreePrimeRunningLcm.lean}{510}{exponent_unique_in_short_interval}\\
first-coordinate projection & \lref{Erdos269/ThreePrimeRunningLcm.lean}{544}{smoothExponentShell_card_le_dropFirst}\\
third-coordinate projection & \lref{Erdos269/ThreePrimeRunningLcm.lean}{586}{smoothExponentShell_card_le_dropThird}\\
sorted quadratic estimate & \lref{Erdos269/ThreePrimeRunningLcm.lean}{630}{sorted_pair_quadratic}\\
quadratic shell bound & \lref{Erdos269/ThreePrimeRunningLcm.lean}{647}{smoothExponentShell_card_quadratic}\\
dyadic internal power & \lref{Erdos269/ThreePrimeRunningLcm.lean}{663}{DyadicInternalPower}\\
internal-power uniqueness & \lref{Erdos269/ThreePrimeRunningLcm.lean}{669}{dyadicInternalPower_exponent_unique}\\
dyadic block base & \lref{Erdos269/ThreePrimeRunningLcm.lean}{751}{dyadicBlockBase235}\\
exact dyadic radix alphabet & \lref{Erdos269/ThreePrimeRunningLcm.lean}{762}{dyadicBlockBase235_cases}\\
bounded-radix consequence & \lref{Erdos269/ThreePrimeRunningLcm.lean}{774}{dyadicBlockBase235_mem_interval}\\
rank-two certificate & \lref{Erdos269/ThreePrimeRunningLcm.lean}{827}{kernel_235_minor_eq_neg_one_fifteen}\\
no integer rotation orbit & \lref{Erdos269/KernelCarryRank.lean}{154}{noIntegerOrbit_logb_of_prime}\\
height factorisation & \lref{Erdos269/KernelCarryRank.lean}{218}{threePrimeHeight_factorisation}\\
kernel factorisation & \lref{Erdos269/KernelCarryRank.lean}{259}{threePrimeKernelQ_factorisation}\\
two-prime outer product & \lref{Erdos269/KernelCarryRank.lean}{278}{twoPrimeKernelQ_eq_outer_product}\\
two-prime vanishing minors & \lref{Erdos269/KernelCarryRank.lean}{298}{twoPrimeKernelQ_minor_two_eq_zero}\\
uniform minors, general form & \lref{Erdos269/KernelCarryRank.lean}{313}{exists_uniform_nonsingular_threePrimeKernel_minor}\\
uniform minors for primes & \lref{Erdos269/KernelCarryRank.lean}{376}{exists_uniform_nonsingular_threePrimeKernel_minor_of_prime}\\
no finite separation & \lref{Erdos269/KernelCarryRank.lean}{388}{not_finite_separable_threePrimeKernel}\\
second-order minor & \lref{Erdos269/KernelCarryRank.lean}{560}{kernel_235_minor2_eq_neg_one_fifteen}\\
third-order minor & \lref{Erdos269/KernelCarryRank.lean}{570}{kernel_235_minor3_eq_one_over_81000}\\
misleading proportional row & \lref{Erdos269/KernelCarryRank.lean}{583}{kernel_235_row_three_eq_smul_row_zero}\\
failure of that proportionality & \lref{Erdos269/KernelCarryRank.lean}{592}{kernel_235_row_three_ne_smul_row_zero_at_four}\\
ordered block digit & \lref{Erdos269/DyadicBlockThresholdPartition.lean}{150}{dyadicOrderedBlockDigit235}\\
normalised state step & \lref{Erdos269/DyadicOrderedTailRecurrence.lean}{110}{dyadicNormalizedTailStateR235_succ}\\
shell summability & \lref{Erdos269/DyadicShellSummability.lean}{142}{summable_dyadicShellMassR235}\\
infinite tail recurrence & \lref{Erdos269/DyadicShellSummability.lean}{177}{dyadicNormalizedShellTsumTailR235_succ}\\
integer-or-cofinally-far dichotomy & \lref{Erdos269/DyadicShellSummability.lean}{188}{dyadicShellTsumTail_integer_or_cofinal_far}\\
boundary divisibility & \lref{Erdos269/RationalLatticeReduction.lean}{51}{smoothHeight_mul_prime_dvd_boundaryHeight}\\
half-height identity & \lref{Erdos269/RationalLatticeReduction.lean}{101}{two_mul_heightNormalizer235}\\
window clearing identity & \lref{Erdos269/RationalLatticeReduction.lean}{158}{heightNormalizer235_mul_windowMass_eq_int}\\
all-scale rationality lattice & \lref{Erdos269/RationalLatticeReduction.lean}{196}{qsmul_normalizedTailState_eq_int_of_value_eq_rat}\\
normalised-state collision & \lref{Erdos269/RationalLatticeReduction.lean}{296}{exists_normalizedTailState_collision_of_value_eq_rat}\\
least positive residue & \lref{Erdos269/ResidueEscape.lean}{26}{leastPositiveResidue}\\
positive representative range & \lref{Erdos269/ResidueEscape.lean}{31}{leastPositiveResidue_pos_le}\\
representative congruence & \lref{Erdos269/ResidueEscape.lean}{52}{leastPositiveResidue_modEq}\\
escape condition & \lref{Erdos269/ResidueEscape.lean}{71}{ResidueEscapesWindow}\\
finite natural-state contradiction & \lref{Erdos269/ResidueEscape.lean}{76}{no_bounded_positive_state_of_residue_escape}\\
contrapositive form & \lref{Erdos269/ResidueEscape.lean}{96}{residue_le_bound_of_bounded_positive_state}\\
integer least-positive-residue obstruction & \lref{Erdos269/ResidueEscape.lean}{110}{no_bounded_positive_int_state_of_leastPositiveResidue}\\
exact residue classifier & \lref{Erdos269/ResidueEscape.lean}{138}{leastPositiveResidue_eq_natAbs_of_pos_le_modEq}\\
window base & \lref{Erdos269/RestrictedFloorSum.lean}{417}{windowBase}\\
window forcing & \lref{Erdos269/RestrictedFloorSum.lean}{422}{windowForcing}\\
affine window identity & \lref{Erdos269/RestrictedFloorSum.lean}{455}{affineRecurrence_window}\\
scaled forcing identity & \lref{Erdos269/RestrictedFloorSum.lean}{469}{windowForcing_const_mul}\\
integral carry window & \lref{Erdos269/RestrictedFloorSum.lean}{480}{integralCarry_window}\\
endpoint residue identification & \lref{Erdos269/RestrictedFloorSum.lean}{497}{leastPositiveResidue_windowForcing_eq_carry}\\
absorbed smooth factor & \lref{Erdos269/RestrictedFloorSum.lean}{548}{smooth3Val_dvd_threePrimeHeight_of_le}\\
common-factor cancellation & \lref{Erdos269/RestrictedFloorSum.lean}{581}{integralCarry_cancel_commonFactor}\\
reduced bound transfer & \lref{Erdos269/RestrictedFloorSum.lean}{601}{reducedCarry_pos_le_of_commonFactor}\\
reduced window transfer & \lref{Erdos269/RestrictedFloorSum.lean}{612}{reducedIntegralCarry_window}\\
cofinal window hypothesis & \lref{Erdos269/RestrictedFloorSum.lean}{629}{CofinalLocalWindowEscape}\\
reduced-carry extinction & \lref{Erdos269/RestrictedFloorSum.lean}{645}{no_positive_reducedCarry_of_cofinalLocalWindowEscape}\\
absorbed-carry extinction & \lref{Erdos269/RestrictedFloorSum.lean}{689}{no_positive_absorbedCarry_of_cofinalLocalWindowEscape}\\
real window identity & \lref{Erdos269/CofinalWindowEscapeEquivalence.lean}{68}{trueNormalizedState_window}\\
escape from irrationality, general bound & \lref{Erdos269/CofinalWindowEscapeEquivalence.lean}{327}{cofinalLocalWindowEscape_of_irrational_of_beaten}\\
escape from irrationality, quadratic family & \lref{Erdos269/CofinalWindowEscapeEquivalence.lean}{355}{cofinalLocalWindowEscape_of_irrational_of_quadratic}\\
escape from irrationality, actual bound & \lref{Erdos269/CofinalWindowEscapeEquivalence.lean}{379}{cofinalLocalWindowEscape_of_irrational}\\
producer equivalence, shifted tail & \lref{Erdos269/CofinalWindowEscapeEquivalence.lean}{392}{actualCofinalLocalWindowEscape_iff}\\
producer equivalence, series value & \lref{Erdos269/CofinalWindowEscapeEquivalence.lean}{398}{actualCofinalLocalWindowEscape_iff_irrational_value}\\
rationality-to-carry bridge & \lref{Erdos269/RationalityCarryBridge.lean}{324}{exists_reducedCarry_of_value_eq_rat}\\
the actual producer & \lref{Erdos269/RationalityCarryBridge.lean}{397}{ActualCofinalLocalWindowEscape}\\
irrationality from the producer & \lref{Erdos269/RationalityCarryBridge.lean}{479}{irrational_value_of_cofinalLocalWindowEscape}\\
bounded-radix alternative & \lref{Erdos269/BoundedRadixTailEscape.lean}{89}{boundedRadix_zero_or_cofinal_far}\\
rational value from an integral scaled tail & \lref{Erdos269/BoundedRadixTailEscape.lean}{183}{rational_of_scaledTail_integer}\\
channel potential classifier & \lref{Erdos269/ThreeChannelBlockRigidity.lean}{59}{channelBlockNull_iff_channelPotential}\\
one-index lift extinction & \lref{Erdos269/CarryLiftExtinction.lean}{178}{no_carryLift_of_errorBound_below_twoPow}\\
uniform lift extinction & \lref{Erdos269/CarryLiftExtinction.lean}{238}{no_bounded_carryLift_of_blockNull_twoAnchors}\\
first-block sum & \lref{Erdos269/CarryLiftExtinction.lean}{289}{binaryLift_twoAnchors_force_firstBlock_sum_one}\\
four-state obstruction & \lref{Erdos269/CarryLiftExtinction.lean}{308}{no_unitAccurateLift_with_twoAnchors_and_firstTwoBlockNull}\\
residue-coboundary form & \lref{Erdos269/WeightedPhaseCarry.lean}{109}{carry_eq_residueDigit_add_coboundary}\\
carry interval & \lref{Erdos269/WeightedPhaseCarry.lean}{150}{carryResidue_mem_interval}\\
digit interval & \lref{Erdos269/WeightedPhaseCarry.lean}{157}{residueDigit_mem_interval}\\
symbolic realisation & \lref{Erdos269/WeightedPhaseCarry.lean}{293}{CartierRealisation}\\
finite realised span & \lref{Erdos269/WeightedPhaseCarry.lean}{334}{finite_realisedSpan_of_factorisation}\\
Cantor-state confinement & \lref{Erdos269/PurePowerIrrationality.lean}{61}{state_mem_Ioo}\\
state gap bound & \lref{Erdos269/PurePowerIrrationality.lean}{74}{stateGap_lower_bound}\\
small-form criterion & \lref{Erdos269/PurePowerIrrationality.lean}{94}{irrational_of_small_nonzero_forms}\\
conditional assembly & \lref{Erdos269/PurePowerIrrationality.lean}{156}{irrational_of_clearing_and_small_gaps}\\
staircase index engine & \lword{Shared/IrrationalRotationStaircase.lean}{271}{exists_staircase_indices}{staircase indices}\\
\end{longtable}
% ---- part family_catalogue ----
\clearpage
\section{Complete result-family map}\label{sec:erdos-269-complete-family-map}

This section places every registered family for this problem in the shared 70-family reader order.  The five display bands control exposition only; the separate promotion state currently covers 6 families and is reported but does not hide or strengthen any family.  Mathematical statements and evidence modes come from the public claim registry.  Across all eight problems the public result-atom catalog contains 681 exact packet coordinates; this problem contributes 65.  The recovered source catalog contains 682 rows; 1 row(s) belong to families absent from the current claim registry and remain disclosed as detached source rows.  Catalog rows expose bounded statement excerpts plus full-source digests, not a claim that every complete packet statement is reproduced here.

\subsection{Three prime lcm cells}
\textbf{Reader position.} 6 of 70; display band: front door.  Formal editorial disposition: merge.  These are separate classifications.

\textbf{Reader entry.} Exact running-LCM product, logarithmic-cell constancy, coordinate jump ratios, and jump count.

Exact running-LCM product, logarithmic-cell constancy, coordinate jump ratios, and jump count.

\textbf{Authority and reach.} Lean kernel; Comparator-selected; locally proved result; novelty unassessed.

\textbf{Exact boundary.} Cell structure alone does not prove irrationality.

\textbf{Result-atom population.} 5 of 681 public coordinates. The public atom catalog groups them under this family.

\textbf{Comparator assurance.} exact selected interface; 1 executable interface(s), 1 repository-registered selected result interface(s).

\textbf{Publication placement.} This family is also admitted to the short note.

\subsection{Two prime transcendence}
\textbf{Reader position.} 17 of 70; display band: major result.  Formal editorial disposition: hold.  These are separate classifications.

\textbf{Reader entry.} Both two-prime running-lcm series are proved transcendental by an authored deduction from an external theorem.

Both two-prime running-lcm series are proved transcendental by an authored deduction from an external theorem.

\textbf{Authority and reach.} paper argument plus cited theorem; paper plus external theorem.

\textbf{Exact boundary.} Comparator cannot certify the external analytic input.

\textbf{Result-atom population.} 1 of 681 public coordinates. The public atom catalog groups them under this family.

\textbf{Comparator assurance.} not applicable to comparator; 0 executable interface(s), 0 repository-registered selected result interface(s).

\textbf{Publication placement.} This family is also admitted to the short note.

\subsection{Weighted phase carry observer}
\textbf{Reader position.} 24 of 70; display band: mechanism.  Formal editorial disposition: hold.  These are separate classifications.

\textbf{Reader entry.} An exact weighted-phase carry recurrence splits into a finite residue digit and an uncontrolled integral quotient coboundary; an explicit function-faithful finite-dimensional observer then forces finite realised span.

An exact weighted-phase carry recurrence splits into a finite residue digit and an uncontrolled integral quotient coboundary; an explicit function-faithful finite-dimensional observer then forces finite realised span.

\textbf{Authority and reach.} Lean kernel; locally proved result; novelty unassessed.

\textbf{Exact boundary.} The recurrence supplies a finite residue coordinate but leaves an uncontrolled integral quotient coboundary; it proves neither a finite-state quotient nor a literal infinite realised span. Finite realised span requires an explicit factorisation through a finite-dimensional function-faithful observer, and scalar evaluation alone is insufficient. No rationality or irrationality conclusion follows; the actual three-prime running-LCM bridge and cofinal escape remain open.

\textbf{Result-atom population.} 7 of 681 public coordinates. The public atom catalog groups them under this family.

\textbf{Comparator assurance.} exact selected interface; 1 executable interface(s), 0 repository-registered selected result interface(s).

\textbf{Publication placement.} This family remains in the complete long record and is not a short-note headline.

\begin{itemize}
  \item \nolinkurl{ErdosProblems.Erdos269.carry_eq_residueDigit_add_coboundary}
  \item \nolinkurl{ErdosProblems.Erdos269.carryResidue_mem_interval}
  \item \nolinkurl{ErdosProblems.Erdos269.residueDigit_mem_interval}
  \item \nolinkurl{ErdosProblems.Erdos269.finite_realisedSpan_of_factorisation}
\end{itemize}

\subsection{Height fibre and shell}
\textbf{Reader position.} 28 of 70; display band: mechanism.  Formal editorial disposition: merge.  These are separate classifications.

\textbf{Reader entry.} Finite height-fibre normal form and a quadratic smooth-shell multiplicity bound.

Finite height-fibre normal form and a quadratic smooth-shell multiplicity bound.

\textbf{Authority and reach.} Lean kernel; locally proved result; novelty unassessed.

\textbf{Exact boundary.} The fibre bounds do not provide the missing divisibility bridge.

\textbf{Result-atom population.} 5 of 681 public coordinates. The public atom catalog groups them under this family.

\textbf{Comparator assurance.} represented by selected interface; 0 executable interface(s), 0 repository-registered selected result interface(s).

\textbf{Publication placement.} This family remains in the complete long record and is not a short-note headline.

\begin{itemize}
  \item \nolinkurl{ErdosProblems.Erdos269.finiteSmoothKernelSum_groupedByHeight}
  \item \nolinkurl{ErdosProblems.Erdos269.smoothExponentShell_card_quadratic}
\end{itemize}

\subsection{Conditional carry escape}
\textbf{Reader position.} 42 of 70; display band: frontier.  Formal editorial disposition: hold.  These are separate classifications.

\textbf{Reader entry.} Under the denominator-dependent cofinal local-window residue-escape predicate, no positive reduced carry can satisfy the exact multiplier recurrence together with its short bound. The load-bearing consumer is no\_positive\_reducedCarry\_of\_cofinalLocalWindowEscape; an absorbed nonzero common-factor carry reduces exactly to that consumer. CofinalLocalWindowEscape, windowBase, windowForcing, leastPositiveResidue, and the absorbed-carry bridge are subordinate finite-window mechanism evidence.

Under the denominator-dependent cofinal local-window residue-escape predicate, no positive reduced carry can satisfy the exact multiplier recurrence together with its short bound. The load-bearing consumer is no\_positive\_reducedCarry\_of\_cofinalLocalWindowEscape; an absorbed nonzero common-factor carry reduces exactly to that consumer. CofinalLocalWindowEscape, windowBase, windowForcing, leastPositiveResidue, and the absorbed-carry bridge are subordinate finite-window mechanism evidence.

\textbf{Authority and reach.} Lean kernel; Comparator-selected; conditional no-go consumer; novelty and significance unassessed.

\textbf{Exact boundary.} The representative assumes b,m : \ensuremath{\mathbb{N}} \ensuremath{\to} \ensuremath{\mathbb{N}}, shortBound : \ensuremath{\mathbb{N}} \ensuremath{\to} \ensuremath{\mathbb{N}} \ensuremath{\to} \ensuremath{\mathbb{N}}, CofinalLocalWindowEscape b m shortBound, and for each positive B coprime to 30 a positive integer-valued d with d(n+1) = (b n : \ensuremath{\mathbb{Z}}) d(n) \ensuremath{-} (B : \ensuremath{\mathbb{Z}}) (m n : \ensuremath{\mathbb{Z}}) and Int.natAbs (d n) \ensuremath{\le} shortBound B n; it then gives False. The absorbed-carry bridge additionally assumes smoothFactor \ensuremath{\ne} 0, c n = smoothFactor * d n, and the corresponding exact absorbed recurrence, then cancels that factor to the same reduced consumer. The cofinal local-window escape producer and the bridge from the actual three-prime running-LCM series or its rationality to this reduced carry remain open. This is not a \#269 endpoint or irrationality proof, and no actual-series identification, novelty, priority, significance, or external-review claim is made. It is one conditional window-carry family, distinct from the finite residue, rank-two, and weighted-phase observer families.

\textbf{Result-atom population.} 37 of 681 public coordinates. The public atom catalog groups them under this family.

\textbf{Comparator assurance.} exact selected interface; 1 executable interface(s), 1 repository-registered selected result interface(s).

\textbf{Publication placement.} This family is also admitted to the short note.

\begin{itemize}
  \item \nolinkurl{ErdosProblems.Erdos269.CofinalLocalWindowEscape}
  \item \nolinkurl{ErdosProblems.Erdos269.no_positive_reducedCarry_of_cofinalLocalWindowEscape}
  \item \nolinkurl{ErdosProblems.Erdos269.windowBase}
  \item \nolinkurl{ErdosProblems.Erdos269.windowForcing}
  \item \nolinkurl{ErdosProblems.Erdos269.leastPositiveResidue}
  \item \nolinkurl{ErdosProblems.Erdos269.no_positive_absorbedCarry_of_cofinalLocalWindowEscape}
\end{itemize}

\subsection{Rank two kernel no go}
\textbf{Reader position.} 54 of 70; display band: frontier.  Formal editorial disposition: hold.  These are separate classifications.

\textbf{Reader entry.} The 2,3,5 kernel is not rank one and its smallest displayed minor equals -1/15.

The 2,3,5 kernel is not rank one and its smallest displayed minor equals -1/15.

\textbf{Authority and reach.} Lean kernel; Comparator-selected; no-go result.

\textbf{Exact boundary.} Failure of rank one does not itself imply irrationality.

\textbf{Result-atom population.} 3 of 681 public coordinates. The public atom catalog groups them under this family.

\textbf{Comparator assurance.} exact selected interface; 1 executable interface(s), 1 repository-registered selected result interface(s).

\textbf{Publication placement.} This family is also admitted to the short note.

\subsection{Dyadic block alphabet}
\textbf{Reader position.} 62 of 70; display band: technical support.  Formal editorial disposition: merge.  These are separate classifications.

\textbf{Reader entry.} The exact dyadic block alphabet is 2, 6, 10, and 30.

The exact dyadic block alphabet is 2, 6, 10, and 30.

\textbf{Authority and reach.} Lean kernel; locally proved result; novelty unassessed.

\textbf{Exact boundary.} The finite alphabet does not supply the needed carry escape.

\textbf{Result-atom population.} 5 of 681 public coordinates. The public atom catalog groups them under this family.

\textbf{Comparator assurance.} represented by selected interface; 0 executable interface(s), 0 repository-registered selected result interface(s).

\textbf{Publication placement.} This family remains in the complete long record and is not a short-note headline.

\begin{itemize}
  \item \nolinkurl{ErdosProblems.Erdos269.dyadicBlockBase235_cases}
\end{itemize}

\subsection{Three prime finite search}
\textbf{Reader position.} 64 of 70; display band: technical support.  Formal editorial disposition: hold.  These are separate classifications.

\textbf{Reader entry.} A local-window checker searches 106666 denominator and start pairs and records small certificates.

A local-window checker searches 106666 denominator and start pairs and records small certificates.

\textbf{Authority and reach.} external exact computation; finite computation.

\textbf{Exact boundary.} Finite search is not a cofinal statement.

\textbf{Result-atom population.} 2 of 681 public coordinates. The public atom catalog groups them under this family.

\textbf{Comparator assurance.} not applicable to comparator; 0 executable interface(s), 0 repository-registered selected result interface(s).

\textbf{Publication placement.} This family remains in the complete long record and is not a short-note headline.
% ---- part back ----
\begin{thebibliography}{10}
\small
\raggedright
\setlength{\itemsep}{0pt}
\bibitem{erdosgraham1980} P.~Erd\H{o}s and R.~L. Graham,
  \href{https://mathweb.ucsd.edu/~ronspubs/80_11_number_theory.pdf}{\emph{Old
  and New Problems and Results in Combinatorial Number Theory}},
  Monogr. Enseign. Math. 28, Geneva, 1980, p.~65.  For a possibly infinite
  prime set $Q$, the page states the infinite-$Q$ irrationality and asks what
  happens for finite $Q$ with more than one element.
\bibitem{erdos1988} P.~Erd\H{o}s, \emph{On the irrationality of certain series:
  problems and results}, in A.~Baker (ed.), \emph{New Advances in Transcendence
  Theory}, Cambridge UP, 1988, pp.~102--109,
  doi:\href{https://doi.org/10.1017/CBO9780511897184.009}{10.1017/CBO9780511897184.009}.
\bibitem{erdos1974letter} P.~Erd\H{o}s,
  \href{https://www.fq.math.ca/Scanned/12-4/letter.pdf}{\emph{Letter to the
  Editor}} (written 1~January 1973), Fibonacci Quart. \textbf{12} (1974),
  no.~4, p.~335.  The letter poses the full series as a conjecture and asserts
  irrationality after retaining only the distinct running-LCM values, for a
  general finite list of primes, without proof.
\bibitem{erdosproblems} T.~F. Bloom,
  \href{https://www.erdosproblems.com/269}{\emph{Erd\H{o}s Problem \#269}},
  \texttt{erdosproblems.com/269}, accessed 28 July 2026 (page displays
  ``last edited 28 December 2025'').  The current record labels the
  finite-support problem open, cites \texttt{[ErGr80, p.~65]} and
  \texttt{[Er88c, p.~106]}, routes to the 1974 letter on p.~335, records the
  infinite-prime and de-duplicated variants, and explicitly describes its
  status as the website owner's present assessment, with no guarantee of
  literature completeness.
\bibitem{formalconjectures269} The Formal Conjectures Authors,
  \href{https://github.com/google-deepmind/formal-conjectures/blob/f776d2f2039351b00737ffcafb9d7d7666e1d9af/FormalConjectures/ErdosProblems/269.lean}
  {\emph{FormalConjectures.ErdosProblems.\texttt{269}}},
  Lean source at commit \texttt{f776d2f}, 2025, accessed 28 July 2026.
\bibitem{apostol1976} T.~M. Apostol,
  \href{https://doi.org/10.1007/978-1-4757-5579-4}
  {\emph{Introduction to Analytic Number Theory}},
  Springer, New York, 1976.
\bibitem{hildebrand1986} A.~Hildebrand,
  \emph{On the number of positive integers $\le x$ and free of prime factors
  $>y$}, J. Number Theory \textbf{22} (1986), 289--307,
  \href{https://doi.org/10.1016/0022-314X(86)90013-2}{DOI}.
\bibitem{montgomeryvaughan2007} H.~L. Montgomery and R.~C. Vaughan,
  \emph{The Prime Number Theorem}, in \emph{Multiplicative Number Theory I:
  Classical Theory}, Cambridge Studies in Advanced Mathematics 97,
  Cambridge University Press, 2007, pp.~168--198,
  doi:\href{https://doi.org/10.1017/CBO9780511618314.008}
  {10.1017/CBO9780511618314.008}.
\bibitem{kovactao2024} V.~Kova\v{c} and T.~Tao,
  \href{https://doi.org/10.1007/s10474-025-01528-0}
  {\emph{On several irrationality problems for Ahmes series}},
  Acta Math. Hungar. \textbf{175} (2025), 572--608,
  doi:\href{https://doi.org/10.1007/s10474-025-01528-0}
  {10.1007/s10474-025-01528-0};
  \href{https://arxiv.org/abs/2406.17593}{arXiv:2406.17593}, 2024.
\bibitem{bugeaudlaurent2023} Y.~Bugeaud and M.~Laurent,
  \emph{Transcendence and continued fraction expansion of values of
  Hecke--Mahler series}, Acta Arith. \textbf{209} (2023), 59--90,
  doi:10.4064/aa220323-18-1;
  \href{https://irma.math.unistra.fr/~bugeaud/travaux/BuMLAA.pdf}
  {authors' Online First PDF},
  \href{https://arxiv.org/abs/2203.12901}{arXiv:2203.12901v1}.
  Theorem~1.1 appears on p.~3 of the linked Online First PDF; its
  introduction, p.~2, attributes the zero-shift case to Loxton and van der
  Poorten. The original result is Theorem~8, p.~40, of~\cite{loxtonvdp1977}.
\bibitem{loxtonvdp1977} J.~H. Loxton and A.~J. van der Poorten,
  \emph{Arithmetic properties of certain functions in several variables III},
  Bull. Austral. Math. Soc. \textbf{16} (1977), 15--47,
  \href{https://doi.org/10.1017/S0004972700022978}{doi:10.1017/S0004972700022978};
  Theorem~8, p.~40.
\bibitem{fan2026comment} S.~Fan, comment on Erd\H{o}s Problem \#269,
  \href{https://www.erdosproblems.com/forum/thread/269}
  {erdosproblems.com forum, thread 269}, 26 June 2026.  The comment gives the
  running-LCM identity for every $|P|$, the two-channel factorisation, the
  Hecke--Mahler reduction, and the transcendence conclusion for $|P|=2$;
  follow-up comments there note the extension to arbitrary coprime pairs.
\end{thebibliography}

\end{document}
